\documentclass[preprint,review,12pt]{elsarticle}

\usepackage[T1]{fontenc}
\usepackage[utf8]{inputenc}
\usepackage{graphicx}
\usepackage{booktabs}
\usepackage{tabularx}
\usepackage{array}
\usepackage{url}
\usepackage{xurl}
\usepackage{textcomp}
\usepackage{float}
\usepackage[margin=1in]{geometry}
\emergencystretch=3em
\sloppy

\journal{Computers \& Security}

\begin{document}

\begin{frontmatter}
\title{Field-Level Discrepancy Typing across Aligned Vulnerability Databases}
\author{TODO Author}
\address{TODO affiliation}
\begin{abstract}
Vulnerability records are often aligned by CVE identifier before they are used for security analytics, but alignment does not make their field values agree. This paper studies post-alignment field-level discrepancies between NVD and GHSA records. From 100,032 normalized NVD records and 28,785 normalized GHSA records, the current corpus contains 8,066 CVE-ID aligned pairs. The central design is type-first: a deterministic baseline assigns each target field to one of five working discrepancy types before any source-support question is considered. The resulting RQ1 distribution shows that disagreement is field-specific: severity and affected\_versions retain visible factual-conflict tails, published/date is dominated by representation and temporal differences, references are mostly incomplete, and cwe\_ids are mostly equivalent. We then report non-gold RQ2 rule-trigger and annotation-readiness diagnostics, plus silver-label RQ3 evidence diagnostics for severity and affected\_versions. These artifacts characterize evidence availability, source provenance, threshold sensitivity, and audit targets; they are not human gold and are not final evaluation claims. The main contribution is therefore a scoped task definition, field-aware discrepancy taxonomy, and reproducible baseline artifact for structured vulnerability-data discrepancy typing, with an evidence-constrained adjudication scaffold whose validation is left to the pending gold-label phase.
\end{abstract}
\begin{keyword}
vulnerability \sep NVD \sep GHSA \sep CVE \sep discrepancy \sep adjudication \sep evidence
\end{keyword}
\end{frontmatter}

\section*{Highlights}
\begin{itemize}
\item Defines post-alignment discrepancy typing for NVD-GHSA vulnerability records.
\item Uses type-first triage before source-support adjudication.
\item Separates representation, incompleteness, timing, and factual conflicts.
\item Reports deterministic field-discrepancy distributions over 8,066 CVE pairs.
\item Characterizes evidence availability, provenance, and audit-risk diagnostics.
\end{itemize}

\section*{Draft Status and Claim Boundaries}
This generated LaTeX source is a packaging scaffold produced from the Markdown manuscript and logged artifacts. RQ1 results are deterministic baseline distributions, RQ2 remains a non-gold rule-trigger diagnostic with blank annotation templates, and RQ3 results are silver-label prototype evaluations rather than human-gold performance claims.

\section{Introduction}

Public vulnerability databases are operational infrastructure for vulnerability management, dependency risk analysis, and security research \cite{NVDUsers2024,Churakova2026VEX,Mandl2026VulnDetection}. A single CVE can appear in multiple sources, and those sources often expose overlapping structured fields: severity, publication date, references, weakness identifiers, and affected-version ranges. CVE-ID alignment is therefore necessary, but it is not sufficient. Once two records are aligned, downstream systems still need to decide whether their field values are equivalent, merely formatted differently, incomplete, temporally shifted, or factually incompatible.

This problem has become more operationally visible as vulnerability publication volume has increased. NIST's 2026 NVD operations update reports changes to NVD enrichment practice under record CVE growth, including prioritization of actively exploited and high-severity vulnerabilities \cite{NVDOperations2026}. That process reality makes cross-source comparison more than a cleanup exercise: downstream users may encounter aligned records whose enriched fields differ because the sources have different curation workflows, schemas, publication paths, and update timing \cite{NVDProcess2026,NVDStatus2026,GHSAAbout2026,GHSARepo2026}.

This paper studies that post-alignment problem for NVD and GHSA vulnerability records. The alignment step is already done: the normalized corpus contains 100,032 NVD records, 28,785 GHSA records, and 8,066 aligned CVE-ID pairs. The remaining problem is not record matching, but how to handle field values that still disagree after alignment.

The current deterministic baseline already shows that disagreement is uneven across fields. After the 2026-07-15 input-integrity refresh, severity contains 1,749 factual conflicts, affected\_versions contains 651 factual conflicts, published contains 1,897 temporal discrepancies, and references is dominated by incompleteness. The cwe\_ids field is mostly equivalent, with a smaller tail of incomplete and factual-conflict cases. The affected\_versions count had previously been reduced from 1,272 to 652 by rule tightening and then from 652 to 651 by filtering NVD CPE matches explicitly marked \texttt{vulnerable=false}. These numbers are baseline outputs, not gold labels.

The central design choice is to separate discrepancy typing from adjudication: triage before truth. Treating every mismatch as a factual conflict would overstate the problem: a same-day timestamp formatting difference is not the same as incompatible severity values, and a one-sided reference list is not the same as a contradiction. The current pipeline therefore uses five working categories: equivalent, representation\_discrepancy, incomplete, temporal\_discrepancy, and factual\_conflict. Only the last category is passed to evidence-constrained adjudication.

The current manuscript answers three research questions at different maturity levels. RQ1 is descriptive: what distribution of discrepancy types appears across aligned NVD-GHSA fields? RQ2 is currently diagnostic: which deterministic typing rules explain the observed binary differences, and what boundaries appear in complete-impact non-human audits of references and CWE candidates, pending a human discrepancy-typing gold set? RQ3 asks what evidence availability and source-support patterns appear for sampled severity and affected\_versions conflict candidates. Severity remains a silver-label prototype; affected\_versions additionally has two prediction-sealed, CVE-disjoint holdouts reviewed by blind Codex agents, but still no human gold.

The contributions are correspondingly scoped. The paper defines the post-alignment structured-field discrepancy task, operationalizes a five-category field taxonomy for NVD-GHSA pairs, reports reproducible baseline distributions and rule-trigger diagnostics, and specifies an abstention-capable evidence scaffold for sampled factual conflicts. The first affected\_versions holdout exposes a protocol failure: evaluating source selection over all agreed rows mixes discrepancy typing with adjudication conditional on a true factual conflict. A second untouched holdout preregisters these as separate endpoints; its highly selective type candidate and branch-based source method do not exceed full-coverage or fixed-source comparators. The implementation therefore contributes an auditable evaluation protocol and failure analysis, not a validated affected\_versions performance improvement. Human gold remains pending throughout.

\section{Background and Problem Definition}

The paper uses CVE IDs as the record-alignment key \cite{CVE2026} and studies field values drawn from NVD \cite{NVDAPI2026} and GHSA \cite{GHSA2026}. NVD enrichment adds analysis beyond the base CVE record, including CVSS metrics, CWE mappings, reference tagging, and vulnerable-product applicability information \cite{NVDProcess2026,NVDStatus2026}. GHSA combines GitHub-reviewed advisories, unreviewed community advisories, and NVD-imported records, and it exposes advisory data through an OSV-compatible JSON representation for supported ecosystems \cite{GHSAAbout2026,GHSARepo2026}. Severity follows CVSS-style score/vector representations \cite{CVSS31}, weakness categories are represented through CWE identifiers \cite{CWE2026}, and affected-version ranges are close to the OSV/GHSA advisory schema family \cite{OSVSchema2026}. These sources are not interchangeable: they differ in curation workflow, schema choices, publication timing, and the amount of advisory context exposed through references.

Affected-version comparison is especially schema-sensitive. OSV separates package ecosystem/name identity from affected version lists and range events, including introduced/fixed event boundaries, and GHSA maintains repository-specific extensions for version-range information that is not always expressible in plain OSV fields \cite{OSVSchema2026,GHSARepo2026}. The affected\_versions field is therefore treated as a structured vulnerability-record field whose schema includes package and range information; the current prototype uses token-support and structural diagnostics rather than semantic version-range adjudication.

Published/date discrepancies are also process-sensitive. NVD enrichment, GHSA review, and NVD-imported advisory flows can expose different timestamps for the same CVE because records move through different publication and update paths \cite{NVDProcess2026,NVDOperations2026,Segal2026GHSA}. The temporal\_discrepancy category is intended to keep those timing differences separate from substantive factual conflicts.

We treat an aligned NVD-GHSA pair as input and a field as the unit of analysis. For an aligned pair \texttt{(r\_NVD, r\_GHSA)} and field \texttt{f}, the first task is to assign a discrepancy type to \texttt{(r\_NVD.f, r\_GHSA.f)}. The second task is activated only for factual conflicts: given field values and fetched advisory evidence, record whether the available evidence cues match NVD, GHSA, both sources, neither source, or whether the case should abstain for audit.

The method uses five working labels: equivalent, representation\_discrepancy, incomplete, temporal\_discrepancy, and factual\_conflict. Equivalent means the normalized values carry the same field meaning. Representation\_discrepancy means the raw strings or schemas differ but the intended fact is treated as compatible by the deterministic rule. Incomplete means one side contains field content that the other side lacks. Temporal\_discrepancy captures publication/date differences that are better explained as timing differences than as contradictory facts. Factual\_conflict is reserved for residual cases where the implemented rules do not find compatibility, incompleteness, or temporal explanation.

This distinction matters because downstream actions differ by type, and prior work shows that vulnerability-data quality, source consistency, and scanner behavior affect operational use \cite{NVDUsers2024,Churakova2026VEX,Mandl2026VulnDetection}. Equivalent and representation discrepancies can usually be normalized away. Incompleteness suggests augmentation rather than source selection. Temporal discrepancies require timeline interpretation. Factual conflicts require evidence and may still end in abstention if the evidence is weak or contradictory.

The current labels are useful for organizing the analysis, but they are not gold truth. The method therefore avoids treating all field mismatches as factual conflicts or collapsing record alignment into the discrepancy task. The working problem is two-stage: first normalize and classify the field disagreement; then, if the field is a factual conflict, use evidence from references and related advisory material to decide whether one side is better supported. The current implementation records this problem definition and prototype; it does not claim that the adjudication stage is final.

\section{Method}

\begin{figure}[!htbp]
\centering
\includegraphics[width=\textwidth]{figures/method_framework.png}
\caption{Method framework for post-alignment NVD-GHSA field discrepancy typing and evidence-constrained adjudication. The adjudication stage is currently a silver-label prototype, not a human-gold source-truth result.}
\label{fig:method-framework}
\end{figure}


The method is an implemented baseline pipeline for post-alignment vulnerability-record fusion. Its input is an aligned pair \texttt{(r\_NVD, r\_GHSA)} keyed by CVE ID, and its unit of analysis is a target field within that pair. The pipeline has three operational stages: field-view construction, deterministic discrepancy typing, and an evidence-constrained source-support scaffold for sampled factual conflicts. The first two stages operate on all 8,066 aligned NVD-GHSA pairs. The third stage is evaluated with non-human annotations: silver-label development samples for severity and affected\_versions, plus two CVE-disjoint, prediction-sealed affected\_versions holdouts reviewed independently by Codex agents. It remains a controlled adjudication scaffold rather than a final source-of-truth system.

The design follows a type-first principle: triage before truth. The first decision asks what kind of disagreement is present. The second decision, activated only for residual factual conflicts, asks which source value has matching fetched-evidence cues, whether both values are provisionally defensible, or whether the case should be abstained for manual review. This separation is essential for this corpus because many binary differences are not contradictions: timestamp formatting, one-sided reference lists, and schema-level version-range differences require different downstream actions from incompatible severity labels.

The innovation claim is therefore deliberately narrow. The contribution is the post-CVE-alignment structured-field discrepancy task, the field-aware taxonomy and auditable baseline implementation, and the abstention-capable evidence scaffold for factual-conflict samples. The core method idea is the routing contract itself: a structured-field mismatch is first assigned an operational discrepancy type, and source-support analysis is delayed until the type indicates a residual factual conflict. It is not a new generic truth-discovery algorithm, not a final source-of-truth resolver, and not a semantic affected\_versions adjudicator.

The routing contract is summarized below. Its purpose is to make the routing decision explicit: the discrepancy type determines the next action before any source-support question is asked. This is the main difference from a winner-first fusion view, where every mismatch is treated as a conflict to be resolved.

\begin{table}[!htbp]
\centering
\setlength{\tabcolsep}{3pt}
\footnotesize
\caption{Method contract for routing field discrepancies before source-support adjudication. Validation boundaries reflect the current validation state.}
\label{tab:03-method-table-1}
\begin{tabularx}{\textwidth}{>{\raggedright\arraybackslash}p{0.18\textwidth}>{\raggedright\arraybackslash}X>{\raggedright\arraybackslash}X>{\raggedright\arraybackslash}X}
\toprule
Discrepancy type & Typical trigger in this study & Downstream action & Current validation boundary \\
\midrule
equivalent & Canonical severity, CWE set, or affected-version span matches & Merge or suppress as non-actionable difference & Deterministic baseline; pending RQ2 gold check \\
representation & Same-day timestamps, label aliases, URL/range-shape formatting & Normalize representation and keep provenance & Deterministic baseline; pending RQ2 gold check \\
incomplete & One-sided references, missing CWE IDs, one-sided affected-version content & Treat as augmentation or coverage gap, not source selection & Deterministic baseline; pending RQ2 gold check \\
temporal & Cross-day published/date differences & Interpret as publication or enrichment timing & Deterministic baseline; pending RQ2 gold check \\
factual conflict & Residual incompatible severity labels or unresolved version structures & Construct evidence context and allow NVD, GHSA, both, neither, or abstain & RQ3 prototype only; human audit pending \\
\bottomrule
\end{tabularx}
\end{table}

\subsection{Operational input/output contract}

The method operates on field instances rather than whole records. A field instance is \path{i = (cve_id, f, r_NVD.f, r_GHSA.f)}, where \texttt{f} is one target field and the two records have already been aligned by CVE ID. The implementation does not overwrite either source record. Each stage appends a new record around the original values so that later analysis can trace a decision back to source values, normalization, rule triggers, evidence URLs, and human-audit fields.

\begin{table}[!htbp]
\centering
\setlength{\tabcolsep}{3pt}
\footnotesize
\caption{Operational input/output contract for the implemented method stages. Validation boundaries reflect the current validation state.}
\label{tab:03-method-operational-input-output-contract-1}
\begin{tabularx}{\textwidth}{>{\raggedright\arraybackslash}p{0.18\textwidth}>{\raggedright\arraybackslash}X>{\raggedright\arraybackslash}X>{\raggedright\arraybackslash}X}
\toprule
Stage output & Required content & Used for & Current boundary \\
\midrule
Field view \texttt{V\_i} & Raw source values, normalized comparands, source provenance, rule-note carrier & RQ1 distribution and RQ2 typing audit & Normalization is baseline logic, not validated semantics \\
Typed instance \texttt{T\_i} & \texttt{V\_i}, discrepancy label, rule identifier or rule note, routing action & RQ1 counts, RQ2 rule-trigger diagnostic, FC sampling & Labels are deterministic outputs until RQ2 gold exists \\
Evidence context \texttt{E\_i} & Candidate URLs, URL provenance, fetch status, fetched snippets, source-class diagnostics & RQ3 evidence availability and audit worklist & Snippets are evidence inputs, not verified page reconstructions \\
Adjudication record \texttt{A\_i} & Source-support decision, matched cues, abstention state, required human-audit fields & RQ3 silver diagnostic and later human-gold evaluation & Silver labels are annotation aids, not gold labels \\
Sealed holdout record \texttt{H\_i} & CVE-disjoint cohort commitment, blind evidence worklist, prediction-only seal, dual-review decision, strict consensus & RQ3 affected\_versions leakage and task-definition audit & Dual-Codex consensus is non-human and selective; human signoff remains pending \\
\bottomrule
\end{tabularx}
\end{table}

Procedure 1 runs over the aligned corpus and is deterministic. For each CVE-aligned pair and target field, it preserves both source values, builds field-specific normalized comparands, applies the ordered rules described below, and emits a typed field instance with a rule note and routing action. Non-FC labels stop at the routing layer: they are not passed to source-support adjudication. Procedure 2 is activated only for sampled factual-conflict candidates. It collects candidate URLs from NVD/GHSA reference contexts, records URL provenance and fetch status, extracts text snippets when available, matches diagnostic field-specific cues, and emits a provisional NVD, GHSA, both, neither, or abstain decision for human audit. The guarded evaluator computes human-gold metrics only after valid final human rows exist.

This contract is also the reproducibility boundary. The paper-facing artifacts can show which field instance was processed, which rule routed it, which evidence was available, and which final human fields remain blank. A claim is allowed only if the corresponding stage output exists and the relevant validation boundary has been satisfied.

\subsection{Field-view construction}

The field-view stage extracts the target fields from each aligned pair and normalizes them into comparable representations. The implemented baseline covers severity, published/date, references, affected\_versions, and cwe\_ids. Each field view preserves the original source values, the normalized comparison representation, and a rule note explaining the path taken by the deterministic comparator. This makes the later diagnostics auditable: a field instance can be traced back to the source records, the normalization step, and the exact rule trigger used by the baseline.

Severity labels are mapped to canonical CVSS-style labels so that source-specific labels such as moderate and medium compare under the same vocabulary. Published/date values are parsed and compared after timestamp normalization, while preserving the distinction between same-day representation differences and cross-day timing differences. References are deduplicated by URL and host-level structure so that subset and formatting effects can be separated from residual conflicts. CWE identifiers are canonicalized as sets. Affected\_versions values are converted into span-like records when the source schema allows it, with the current implementation retaining range-bound shape, endpoint strings, and package/context information for later diagnostics.

For references, the development workflow treats URL normalization as a resource-identity hypothesis rather than a string-cleaning fact. A complete-impact audit groups only the URL identities changed by a candidate profile, freezes live-HTTP metadata and content fingerprints, masks the candidate transformation and expected label from two reviewers, and retains a change only under strict two-run agreement. The audited development profile keeps exact HTTP redirects, scoped Liferay presentation-query removal, and identifier-preserving GitHub-advisory and Huntr aliases, but excludes encoded GitHub line-suffix stripping after the reviewers apply incompatible underlying-file and HTTP-resource interpretations. This profile is selected after the audit and is therefore a post-hoc development artifact, not a confirmed normalization rule.

The field-view stage is intentionally conservative. It does not try to infer missing package identity or repair upstream record alignment. As a diagnostic baseline, when affected\_versions values use incompatible package names, ecosystem conventions, or range schemas, the current implementation records those structures for audit rather than silently forcing them into a shared range model.

Validation status: field-view construction is implemented and reproducible over the aligned corpus. Its normalization choices are still baseline rules and should be evaluated through the pending RQ2 human labels before being treated as validated field semantics.

\subsection{Discrepancy typing}

The discrepancy-typing stage applies field-specific deterministic rules and emits one of the five working labels: equivalent, representation\_discrepancy, incomplete, temporal\_discrepancy, or factual\_conflict. Equivalent means the normalized values carry the same field meaning. Representation\_discrepancy means the raw strings or schemas differ but the implemented rule treats the intended field fact as compatible. Incomplete means one side has field content that the other side lacks. Temporal\_discrepancy captures publication/date differences that are better modeled as process timing than as a contradictory value. Factual\_conflict is the residual category used when the implemented rules do not find equivalence, representation compatibility, incompleteness, or temporal explanation.

The rules are field-specific because the same surface operation has different meaning across fields. The general priority is conservative: preserve missingness first when one side lacks comparable content; test field-specific equivalence and representation compatibility; route one-sided coverage gaps to incompleteness; route publication timing differences to temporal\_discrepancy; and use factual\_conflict only as the residual category. When multiple non-conflict explanations are possible, the implementation emits the most operationally specific routing label: temporal for date timing, incomplete for one-sided content, representation\_discrepancy for compatible form/schema differences, and equivalent for normalized identity.

\begin{table}[!htbp]
\centering
\setlength{\tabcolsep}{3pt}
\footnotesize
\caption{Ordered field-specific baseline rule paths for discrepancy typing. The rules are diagnostic until RQ2 human labels exist.}
\label{tab:03-method-discrepancy-typing-1}
\begin{tabularx}{\textwidth}{>{\raggedright\arraybackslash}p{0.18\textwidth}>{\raggedright\arraybackslash}X>{\raggedright\arraybackslash}X}
\toprule
Field & Ordered baseline rule path & Boundary or known limitation \\
\midrule
severity & Canonicalize CVSS-style labels; matching canonical labels route away from conflict; missing one-sided values route to incomplete; different canonical labels become factual-conflict candidates. & Label canonicalization is deterministic; correctness of a source rating is not decided here. \\
published/date & Parse timestamps; same-day differences route to representation\_discrepancy; cross-day differences route to temporal\_discrepancy. & Date routing explains process timing, not which publication timestamp is more authoritative. \\
references & Normalize URLs and host-level structure; URL-format variants route to representation\_discrepancy; strict one-sided or subset reference sets route to incomplete; rare residual disagreement remains factual\_conflict. & Reference provenance is not evidence correctness; a URL's source side does not decide truth. \\
cwe\_ids & Canonicalize CWE IDs as sets; set equality routes to equivalent; one-sided or subset content routes to incomplete; residual set disagreement remains factual\_conflict. & CWE granularity and upstream mapping policy are not adjudicated by this baseline. \\
affected\_versions & Preserve package/context, endpoint strings, range-bound shape, and span-like records; baseline structural or effective-bound heuristics can route cases away from conflict; one-sided content routes to incomplete; residual incompatible version structures remain factual-conflict candidates. & Current logic is structural and diagnostic; it is not package-aware or range-aware semantic adjudication. \\
\bottomrule
\end{tabularx}
\end{table}

Examples are used only to clarify diagnostic-only routing behavior. A severity pair such as \texttt{HIGH} versus \texttt{CRITICAL} remains a factual-conflict candidate; a same-day timestamp string difference is a representation discrepancy; cross-day publication values are temporal; a strict-subset reference list is incomplete; and an NVD point version inside a parseable GHSA range is routed as baseline range-shape compatibility, pending human check. For affected\_versions, this remains an audit-prioritization heuristic and diagnostic token-support prototype: package names that do not overlap and version tokens that do not align are kept as audit-risk residuals rather than treated as semantic version-range adjudication.

Affected\_versions requires the most structural handling. The implemented baseline compares span-like records and applies range-bound diagnostics such as effective-bound equivalence, one-sided presence, shared upper-bound endpoints, and the 2026-04-22 rule tightenings that lowered affected\_versions factual-conflict candidates from 1,272 to 652. The 2026-07-15 input-integrity refresh then filtered NVD CPE matches explicitly marked \texttt{vulnerable=false}, lowering the current residual count from 652 to 651. These changes are implementation diagnostics: they show how many cases the pipeline routes differently, not that correctness improved by the same amount. The tightenings include cases where \texttt{end\_including} and \path{end_excluding/fixed} share a major.minor prefix, where an NVD point version falls inside a parseable GHSA range, where \texttt{end\_excluding} strings differ only by prefix truncation, and where NVD uses a date string as a version-end boundary. These rules are audit-prioritization heuristics, not human-validated truth. They can fail when the two sources name different packages, use ecosystem-specific ordering, encode branch- or distribution-specific backports, or express the same advisory through incompatible range schemas. Because the planned 300-instance discrepancy-typing gold set is not complete, the current RQ2 artifact reports rule-trigger distributions and audit targets rather than typing accuracy, precision, recall, or macro-F1.

The diagnostic input-integrity analysis is separately recorded at \path{results/input_integrity/affected_versions/affected_versions_input_integrity.{json,md\}}. The legacy aligned corpus contained 1,105 \texttt{vulnerable=false} CPE items across 106 of the 8,066 matched records. Removing them changed 10 affected\_versions baseline classifications. Seven rows in the frozen 100-sample affected\_versions cohort had refreshed source inputs, but the diagnostic cohort mapping, its 40 determinate versus 60 abstained AI-adjudicated decisions, and all final labels remained unchanged after targeted recheck. This recheck is AI provenance (\texttt{label\_is\_human=false}), not human validation. A scan of the current 28,785-record reviewed GHSA snapshot found no range containing multiple events, so the present first-event flattening did not alter this snapshot; it remains a latent risk for future snapshots rather than a validated general assumption.

Validation status: discrepancy typing is implemented and useful for audit prioritization. It is not yet a gold-backed classifier, so all RQ2 accuracy and agreement claims remain pending.

\subsection{Evidence context construction}

The evidence stage is activated only for sampled factual-conflict candidates. It collects candidate URLs from the aligned records' reference contexts, fetches text snippets and basic metadata, and stores a per-sample evidence context. The evidence context keeps source provenance separate from source truth: a URL can originate from NVD references only, GHSA references only, or both, but that provenance alone does not decide which field value is correct.

For the severity seed, the project assembled 80 examples, collected 470 candidate evidence URLs, and cached fetched evidence text. Of those URL records, 415 returned HTTP-ok fetched text snippets, 46 hit URL errors, 4 returned 403, 3 returned 404, and 2 timed out. All 80 samples have at least one successful text evidence source. For affected\_versions, the same evidence builder has been run on the 100-sample manual-check set. Its manifest records 585 candidate evidence URLs: 402 HTTP-ok responses, of which 401 yielded OK text records after one skipped non-text URL, plus 132 URL errors, 30 timeouts, 15 HTTP 403 responses, 3 HTTP 404 responses, and 2 HTTP 308 responses. All 100 samples have at least one successful text evidence source.

The evidence-source reliability diagnostic further classifies each fetched URL by source class and reference provenance. Source classes distinguish NVD pages, GitHub advisories, GitHub commits or repositories, vendor advisories, package registries, mailing lists or issue trackers, and other pages. These classes are heuristic audit features, not source-truth labels. They are used to characterize evidence availability, identify fetch bias, and prioritize manual review when a field's evidence comes mainly from pages that are hard to fetch or weakly structured.

Validation status: evidence-context construction is implemented for the current severity and affected\_versions samples. The fetched snippets are evidence inputs, not manually verified page reconstructions, and fetchability differences remain a validity risk.

\subsection{Evidence-constrained source-support prototype}

The adjudication decision space is source-support oriented rather than winner-only. For human-audited evaluation, a sampled factual conflict can be labeled as NVD-supported, GHSA-supported, both-source supported, neither-source supported, or abstain. NVD-supported and GHSA-supported mean that the available evidence is judged to support one side's field value. Both-source supported means that both values are defensible or supported under different source contexts. Neither-source supported means that the fetched evidence does not support either recorded value. Abstain means that the evidence is missing, weak, contradictory, fetch-biased, or outside the prototype's cue model. The prototype output therefore records the predicted source decision, matched evidence cues, and an abstention path rather than forcing a value in every case.

The current evidence-aware silver-v2 label sets are generated from the fetched evidence contexts and record provisional source decisions for each sample, including both-source support, NVD-only, GHSA-only, neither-source, and abstain outcomes. Silver-v2 is an annotation aid and diagnostic comparison target, not an evaluation oracle. The severity prototype uses an \texttt{evidence\_score\_baseline} that matches severity labels, numeric scores, and CVSS vectors in fetched text snippets. The affected\_versions prototype uses a \path{version_token_support_baseline} that checks whether fetched text contains version tokens from either side. The latter is intentionally shallow: it does not parse package identity, ecosystem-specific version ordering, or version-range containment, which are known schema-sensitive parts of OSV/GHSA affected-version representation and affected-version identification \cite{OSVSchema2026,GHSARepo2026,Chen2025AffectedVersions}.

Fixed-source baselines (\texttt{prefer\_nvd}, \texttt{prefer\_ghsa}) and a \texttt{latest\_published} baseline provide simple comparisons. Within the silver-label setup, the evidence-aware baselines test whether lightweight evidence matching changes the diagnostic pattern relative to fixed-source preferences, while preserving the distinction between supported, unsupported, and abstained cases. The current results should therefore be read as selective baselines against silver labels, not as final adjudication performance.

For affected\_versions, the v1 holdout protocol excludes every CVE in the development cohort before fixed-hash sampling, uses a dedicated evidence cache, strips prior labels and predictions from the reviewer worklist, and seals label-free method predictions before review. Two Codex reviewers independently assign both a discrepancy type and source support. The merger accepts only determinate exact joint agreement that satisfies the evidence-support contract. Its preregistered endpoint evaluates source selection on all accepted rows; a post-unseal diagnostic then separates true factual conflicts from non-conflicts after discovering that the original endpoint conflates the two tasks.

The v2 protocol turns that failure into an explicit design constraint. It excludes both the development and v1 CVEs, preregisters discrepancy typing and source adjudication on strict factual conflicts as separate endpoints, and evaluates an independent source head without gating it on the method's own type prediction. Reviewer evidence citations are bound to literal frozen snippets through structured endpoint, target, and role claims. Prediction files, method code, protocol code, reviewer identities, row order, and evidence hashes are sealed or checked before evaluation. This protocol prevents the v1 metric conflation; it does not by itself validate either predictor.

Validation status: the adjudication stage is a prototype. Its current role before human audit is to define the source-support decision space, exercise evidence availability, and generate a worklist for manual adjudication.

\subsection{Audit interface and current boundary}

The implementation already prepares the handoff from prototype adjudication to human audit. The RQ2 annotation seed and 20\% consistency-review template cover the target fields and rule-trigger strata, but their labels remain blank. The RQ3 gold-audit templates preserve source identifiers, fetched evidence, silver-v2 decisions, and required human-audit fields for 80 severity rows and 100 affected\_versions development rows, but all final human audit rows are still draft. The guarded RQ3 evaluator refuses to write gold-audit metrics until valid final human rows exist. The affected\_versions holdout consensus rows are also marked \texttt{label\_is\_human=false} and cannot substitute for that audit.

The method stages map to the research questions with different maturity levels. RQ1 uses Procedure 1 over all aligned pairs and target fields to report a deterministic discrepancy distribution. RQ2 uses the same typed instances plus annotation templates and supports rule-trigger/readiness diagnostics together with bounded non-human full-impact audits for references and CWE candidates. RQ3 uses Procedure 2 over sampled factual-conflict candidates and currently supports evidence availability, non-human label diagnostics, holdout agreement/coverage, and evaluation-protocol findings rather than human-gold adjudication performance.

This boundary is part of the method rather than a footnote. The current paper can describe the field-view construction, deterministic typing rules, bounded non-human full-impact audits, evidence-context builder, sealed holdout protocol, abstention mechanism, and audit interface. It cannot yet claim human-gold RQ2 typing performance, validated affected\_versions false-positive rates, an affected\_versions method improvement, or final RQ3 source-truth adjudication performance.

\section{Experimental Setup}

All results in this paper come from reproducible artifacts, not from a human gold set. The main corpus consists of 100,032 normalized NVD records, 28,785 normalized GHSA records, and 8,066 aligned pairs. The current baseline discrepancy view is the 2026-07-15 input-integrity refresh of \texttt{build\_field\_discrepancies.py}.

The claim boundary matrix below records how each artifact is allowed to be used in the manuscript. It is part of the experimental setup because the current package deliberately mixes deterministic baselines, blank annotation templates, and silver labels without treating any of them as human gold.

\begin{table}[!htbp]
\centering
\setlength{\tabcolsep}{3pt}
\scriptsize
\caption{Generated table from 04\_experim.}
\label{tab:04-experim-table-1}
\begin{tabularx}{\textwidth}{>{\raggedright\arraybackslash}p{0.18\textwidth}>{\raggedright\arraybackslash}X>{\raggedright\arraybackslash}X>{\raggedright\arraybackslash}X>{\raggedright\arraybackslash}X}
\toprule
RQ & Artifact & Label source & Allowed claim & Not claimed \\
\midrule
RQ1 & discrepancy stats; RQ1 coverage summary & deterministic field rules & baseline discrepancy distribution across aligned NVD-GHSA fields & gold discrepancy labels or source correctness \\
RQ2 & typing diagnostics; blank annotation templates; references/CWE full-impact audits; fresh-CVE typing-stability cohort & deterministic rules, non-human Codex candidates, strict same-model dual-pass consensus, and blank human fields & rule-trigger/readiness diagnostics, bounded complete-impact candidate findings, and current-snapshot baseline stability & human-gold typing performance, independent-human agreement, candidate-profile gain, or production-rule validation \\
RQ3 severity & severity silver-v2 evidence and predictions & evidence-aware LLM silver labels & silver-label evidence-availability and baseline diagnostics & human-gold adjudication performance or final source truth \\
RQ3 affected\_versions & affected\_versions silver-v2 diagnostics; two CVE-disjoint frozen holdouts & evidence-aware LLM silver labels; strict agreement between blind Codex reviewers & development diagnostics, holdout agreement/coverage, task-separated method diagnostics, and protocol failure analysis & human-gold performance, independent-human validation, or semantic version-range adjudication \\
\bottomrule
\end{tabularx}
\end{table}

The field-coverage table separates raw NVD nonempty coverage from coverage inside the 8,066 aligned NVD-GHSA pairs. This table is descriptive: it explains which fields have one-sided missingness before discrepancy typing, not whether either source is correct.

\begin{table}[!htbp]
\centering
\setlength{\tabcolsep}{3pt}
\scriptsize
\caption{Generated table from 04\_experim.}
\label{tab:04-experim-table-2}
\begin{tabularx}{\textwidth}{>{\raggedright\arraybackslash}p{0.15\textwidth}>{\raggedright\arraybackslash}X>{\raggedright\arraybackslash}X>{\raggedright\arraybackslash}X>{\raggedright\arraybackslash}X>{\raggedright\arraybackslash}X>{\raggedright\arraybackslash}X}
\toprule
Field & NVD nonempty & GHSA nonempty & Both nonempty & NVD only & GHSA only & Both empty \\
\midrule
severity & 95549 & 8066 & 8033 & 0 & 33 & 0 \\
published/date & 100032 & 8066 & 8066 & 0 & 0 & 0 \\
references & 97988 & 8066 & 8037 & 0 & 29 & 0 \\
affected\_versions & 68784 & 8024 & 5667 & 27 & 2357 & 15 \\
cwe\_ids & 90738 & 7949 & 7870 & 60 & 79 & 57 \\
\bottomrule
\end{tabularx}
\end{table}

For RQ2, the analysis uses \path{results/rq2_discrepancy_typing/rq2_typing_diagnostics.json} as a non-gold diagnostic artifact. It aggregates deterministic \texttt{(field, status, rule note)} triggers over the 40,330 field instances in the 8,066 aligned pairs. The repository also contains a 300-instance stratified annotation seed at \path{data/annotations/rq2/discrepancy_typing_seed.jsonl} and \texttt{.csv}, with 60 sampled field instances per target field. A 20\% second-pass consistency review template is prepared under \path{data/annotations/rq2/consistency_review/}, with 60 rows and 12 rows per target field. The readiness diagnostic at \path{results/rq2_discrepancy_typing/rq2_sample_coverage.json} checks that all 17 nonzero field/status candidate strata and the top 12 full-corpus deterministic rule triggers are represented in the primary seed. These files are templates with blank primary and reviewer labels; they are not used to report human-gold accuracy, macro-F1, precision, recall, agreement, or Cohen's kappa.

A separate current-snapshot typing-stability cohort excludes 717 CVEs previously exposed in RQ2/RQ3 development, audits, or holdouts. It samples 250 rows per field with 1,250 globally unique CVEs, projects raw aligned values into inverse-order blind worklists, and seals six executable prediction profiles before either accepted reviewer output. Two Codex CLI pass roles use the same model/config but disjoint ephemeral session sets; strict consensus requires exact determinate labels, no low confidence, and no human-review request. The passes comprise 28 and 67 sessions rather than two persistent annotators. This design estimates non-human baseline stability only: all six sealed prediction columns are identical after exposure exclusion, so candidate-profile gain is structurally unidentifiable on this cohort.

To test whether a post-profile comparison is available, we later acquire an isolated official 2026 NVD feed and a commit-pinned reviewed-GHSA archive after the profile seal, without replacing the main corpus. Eligibility is frozen before review. The strict event-time tier requires both normalized publication timestamps to be later than the seal; only this tier may support a temporal-generalization claim. If it is unavailable, a separately named snapshot-external development tier may use previously unseen \texttt{CVE-2026-*} records under an availability-adaptive size rule. The acquisition finds no strict-event-time CVE but 5,948 snapshot-external single-GHSA CVEs, fixing the latter tier at 50 rows per field. Selection uses a 70\% proportional and 30\% equal-audit allocation over current-status strata, a fixed hash rank, and global CVE uniqueness. It does not use candidate predictions. Six profile columns, inverse-order blind worklists, source hashes, the Codex executable hash, and two pass identifiers are sealed before review. This design can make naturally occurring profile differences identifiable, but it remains development-only because collection time after the seal is not event time after the seal.

After the main comparison is unsealed, a post-selection evidence-secondary contract targets exactly the three rows where current and CWE-taxonomy profiles differ. Before two new reviewer files exist, it freezes nine URLs already listed by NVD or GHSA, their response bodies and metadata, the prior source and evaluation hashes, two inverse-order blind worklists, and literal allowed CWE path strings. Neither worklist contains a profile prediction or prior label. One reviewer judges the taxonomy relation before the concrete mechanism, while the other uses the reverse reasoning order. Strict consensus requires agreement on set relation, discrepancy label, taxonomy support, and whether the mapped weakness is supported by the frozen mechanism evidence. Two earlier protocol attempts are retained but excluded: v1 violates the confidence/uncertainty contract, and v2 violates the literal path interface. V3 uses fresh sessions and validates path values before writing. Because row selection follows the main comparison and all reviewers use the same model family, this stage can explain selected cases but cannot update the sealed 250-row evaluation, estimate method gain, or authorize promotion.

We then broaden the evidence view to all 50 sealed \texttt{cwe\_ids} rows, hiding the three known profile differences among 47 controls. The field-complete packet contains 33 exact sets, nine literal strict subsets, four overlap-non-subsets, and four disjoint pairs. It freezes at most three ranked NVD/GHSA references per row: 134 of 135 URL snapshots succeed. Two reverse-order worklists omit profile predictions, A/B labels, strict status, and profile-difference membership. V1 is rejected because it incorrectly forces every literal strict subset to \texttt{incomplete}; V2 permits evidence-dependent subset judgments but reaches merge with a nonliteral evidence quote. V3 moves rationale length, URL membership, literal-quote, duplicate-citation, and allowed-path checks before output is written. Its E/F roles use 20 disjoint ephemeral sessions. This design removes explicit difference-row selection from the reviewer input, but the CWE field and evidence contract are still chosen after A/B and profile unsealing. The result is therefore a post-hoc field-control diagnostic, not a replacement for the sealed main comparison.

One final evidence-secondary stage covers all 16 non-CWE rows among the 19 original A/B non-strict rows; the other three are excluded because the all-50 CWE audit already gives them strict field-complete decisions. Before retrieval, the stage fixes field counts (12 affected\_versions, two references, and two severity), a deterministic selector over at most six original NVD/GHSA URLs, reverse-order blind G/H worklists, citation gates, and development thresholds of 0.75 evidence availability, 0.40 strict secondary resolution, and 0.95 staged coverage. G/H see neither A/B or CWE decisions nor any profile prediction. The staged candidate may retain the 231 sealed strict rows and add only strict decisions from the two post-selected evidence stages. It is expressly diagnostic: neither passing nor failing may update the sealed profile evaluation, create human gold, or authorize promotion.

After A/B unsealing, a separate tiebreak protocol selects only the 103 non-strict rows from the original blind worklist. Reviewer C receives neither baseline predictions nor A/B outputs, uses the original prompt, and runs in ephemeral sessions disjoint from A/B. Before C output, the tiebreak manifest fixes a qualified-vote rule: at least two determinate, non-low-confidence votes without human-review requests must choose the same label. It also fixes advancement thresholds of 0.70 selected-row resolution and 0.975 combined cohort coverage. Because C shares the same model family, prompt, snapshot, and post-unsealing row selection, its output remains a non-human candidate diagnostic rather than independent-human adjudication.

After that gate fails, a second contract is fixed before new URL retrieval for all 37 remaining non-human candidate abstentions. It selects at most six URLs already present in each original NVD/GHSA reference context, freezes their extracted text/status records, and gives the same evidence-enriched row to reviewers D and E in opposite deterministic orders. D/E receive no baseline, A/B/C output, vote group, or candidate label. For affected\_versions, cwe\_ids, and references, secondary resolution requires strict D/E consensus and at least one successful frozen citation from each reviewer; severity may use its structured CVSS values. The pre-fetch gate requires 0.75 evidence availability, 0.40 secondary strict resolution, and 0.982 combined candidate coverage. This remains a post-unsealing, same-model-family, non-human development audit.

For real-human validation, the repository projects all 1,250 frozen rows into a separate source-bound three-stage packet. Review rows omit baseline status and notes, sampling strata, both Codex decisions, and their consensus. Two different real people must label each row independently before an author records and signs the resolution. An author-only scheduler prioritizes the 103 non-human reviewer disagreements and 216 baseline-versus-consensus disagreements but contains no labels; complete evaluation still requires all 1,250 rows. The validator binds source hashes and rejects source drift, partial pending content, identical reviewer IDs, unsigned resolutions, and any attempt to mark the packet itself as human gold. At the current snapshot all 1,250 rows are pending and none is signed.

The later 250-row post-profile cohort has an analogous full-cohort packet rather than a three-case form. It omits all six sealed predictions, A/B and E/F outputs, consensus, strata, and scheduler signals. Its author-only no-label queue places the 19 original non-strict rows first, then the three profile-difference diagnostics, 44 baseline-versus-consensus mismatches, and 184 completion rows. The blank packet passes ordinary integrity validation but fails both \texttt{--require-signed} and \texttt{--require-complete}; all 250 rows remain pending. File validation cannot establish that an identifier belongs to a real person or that the reviews were independently conducted, so external identity verification and separate human-gold promotion remain mandatory.

To isolate what human labels can change in the profile comparison, a label-free paired-outcome contract reads only the original sealed source and prediction files. It rejects hash drift, identifies all current-versus-CWE prediction differences, and enumerates every five-label assignment on those rows. Each assignment contributes a paired win, loss, or tie according to which sealed prediction matches. The assignment counts are logical cases, not probabilities or a prior over human decisions. A second label-free contract groups all six profiles by their complete prediction vectors and evaluates every profile pair with the conditional exact two-sided McNemar test at alpha 0.05. For the representative cross-vector pair, it enumerates whether each possible label makes current alone correct, the candidate alone correct, or both wrong. It also reports assumption-bound design sensitivities: the theoretical minimum correctness-discordant count for any rejection, conditional exact power under specified directional win probabilities, and future cohort sizes under a stationary observed prediction-difference rate. These diagnostics bound the possible full-cohort accuracy difference and the current cohort's paired-test capacity; they cannot establish absolute accuracy, macro-F1, construct validity, future disagreement prevalence, or the error distribution over equal-prediction rows, which remain part of the full real-person workflow.

Because the 250-row allocation is stratified by current status rather than sampled as a simple random draw, we separately run a prediction-only census over all five fields of all 5,948 snapshot-external eligible CVEs. The census binds the acquisition and profile code, exactly replays every sealed prediction on the selected 250 rows, and only then extends the same six deterministic profiles to 29,740 field instances. It creates neither a replacement cohort nor a reviewer worklist and reads no label, review, consensus, or evidence-secondary output. Its purpose is to measure deterministic profile disagreement availability in this revealed universe and to test whether the sampled 3/250 difference count can be used as a population-rate assumption. It cannot measure correctness, select a winning profile, validate a power model, or restore event-time eligibility.

After that census is revealed, two field-specific audits cover its complete 34-row prediction-difference union without treating the rows as a representative sample. The CWE audit retains all 29 \texttt{cwe\_ids} differences, freezes 85 ranked source-reference records (83 successful), and presents inverse-order worklists with current/candidate values removed. Every row visibly has a disjoint taxonomy relation and an official path, so selection blinding is explicitly incomplete. Strict E/F consensus requires exact agreement on relation, discrepancy label, taxonomy compatibility, and concrete mapping, plus literal frozen evidence. The references audit instead constructs profile-independent partitions over all raw URLs in the five-row references union. It restores NVD/GHSA membership only after review and reports two definitions separately: persistent underlying resource and frozen HTTP resource. Direct side and CVE fields are hidden, but URL text can expose CVE identity and alias structure. The first references seal is archived because its merge calls a nonexistent writer and produces no consensus; revision 2 re-seals the archived attempt, code, 26 URL probes, and fresh reviewer outputs before merge. Both audits are outcome-selected, same-snapshot, same-model-family diagnostics and cannot update the sealed 250-row evaluation or authorize promotion.

Because failure-mode inspection exposed a possible field-contract change, we additionally run a post-hoc cross-protocol stability gate before implementing either mechanism as a comparator. For severity, the fresh missing-score/vector projection is replayed on the 60 older primary rows and 12 older review rows using only their original structured field context. For affected\_versions, the audit compares input availability rather than inventing historical raw values: older seed rows can contain two empty projected range values while retaining package identity on one side, whereas fresh rows expose raw affected records. The advancement gate rejects a candidate if severity direction reverses across protocol generations, affected inputs are not comparable, or no signed human calibration establishes the intended construct. This gate can reject an unstable candidate but cannot validate either protocol.

We then run two sealed development calibrations without changing the comparator. Calibration v1 deterministically selects 60 fresh-holdout rows across four severity boundary patterns, one-sided unbounded affected claims, and unchanged controls. The two pass roles receive identical raw values and one prompt in reverse order through disjoint ephemeral sessions; exact/strict thresholds and per-stratum expected-label thresholds are fixed before review. After v1 exposes four residual constructs, calibration v2 freezes a refined prompt and gate before selecting 42 rows disjoint from every v1 row. V2 separates same-version from cross-version CVSS vectors, equal normalized ranges with different package names, singleton-versus-interval cases, and prerelease boundaries. Open affected-version strata are evaluated for strict reviewer stability rather than forced to share one expected label. For v2's sole non-strict row, a secondary packet freezes structured responses from the official GHSA advisory, fixing commit, and three linked XWiki Jira issues; two new disjoint sessions receive the same evidence, and the secondary gate additionally requires each reviewer to cite at least two frozen official URLs. After that conditional review disagrees, a deterministic projection audit checks component membership, product-to-artifact dependency edges, release-version policy, current and legacy catalog membership, source continuity across artifact migration, and prerelease ordering. Its first pass abstains before historical dependency edges are available; a second post-unsealing pass permits set comparison only after five product dependencies and the legacy/current source transition are frozen. We subsequently fix a generic graph schema and raw-data selector before fetching cross-case evidence. The selector does not read reviewer files: from the already unsealed v2 source it takes all eight two-sided affected\_versions rows with one subject per source, differing full artifact identifiers, and equal non-empty interval signatures. Each row must bind every endpoint using official boundary POMs, root Composer manifests, Go module manifests, or an explicit project/registry alias; missing evidence forces abstention. A second selector keeps the schema unchanged, excludes the XWiki development case, and takes all five remaining one-subject cross-artifact rows with non-equal signatures. Registry catalogs enumerate released versions before applying the fixed equal/subset/overlap relation map. Before running, advancement requires at least 4/5 projected rows and at least 4/5 candidates matching both sealed AI reviewers; reviewer files are read only after graph candidates are computed. After this no-go is observed, we freeze a separate Codex expert-contract candidate that chooses snapshot-extensional published-release semantics without claiming human approval. A label-independent one-to-many selector finds the sole InLong product row with two GHSA Maven components. Each component must bind through its official catalog and same-version InLong parent POM at all three boundaries; affected sets are evaluated separately, their heterogeneity is retained, and their union is compared with the NVD product set only when every edge and affected release is bound. Reviewer labels are consulted only after this candidate is computed. Both calibrations and all conditional post-unsealing audits remain non-human development evidence, not a confirmatory holdout or human gold.

The final graph stress test leaves the reviewed calibration source and selects from the complete aligned input. For each of NuGet, PyPI, and crates.io, it takes the minimum SHA-256-ranked CVE with one NVD product, exactly two GHSA packages, differing canonical component signatures, no more than three spans per claim, and parseable nonzero boundaries. This produces Oracle Database Server, LangChain, and Deno rows without consulting any reviewer or consensus artifact. Before registry retrieval, we fix edge classes for coordinated components, required dependencies, dependency constraints, parallel distributions, and unbound alternatives. A product-level union requires a deterministic total component-release-to-product-release map; an admissible dependency range is insufficient. Advancement requires at least 2/3 complete projections across at least two ecosystems. Fourteen official registry, version-metadata, dependency, and vendor responses are then frozen. This audit remains post-unsealing and non-human, but unlike the calibration-derived graph rows its upstream source is not conditioned on non-human consensus.

After this 0/3 result, we freeze a post-no-go recovery contract before retrieving any additional Deno evidence. The product domain is every stable official \path{denoland/deno} GitHub release from \texttt{1.41.3} through \texttt{2.3.2}, plus its immediate predecessor and successor. Every selected tag must contain one parseable exact \texttt{deno\_runtime} entry in its committed \texttt{Cargo.lock}; that runtime must exist in the frozen crates.io catalog, the complete sequence must be monotonic, and the two outer mappings must bound the claim \path{0.150.0 <= deno_runtime < 0.212.0}. Only then may the GHSA runtime interval be projected into product releases and unioned with the direct Deno intervals. An independent verifier reconstructs the dynamic source inventory and all sets from cached official responses. The fixed advancement threshold is 1/1, but passing remains a single-row, non-human development result.

We separately audit the three non-affected-version rows that remain unresolved after D/E. This cohort is explicitly selected from revealed outcomes, and protocol discovery inspects its evidence shape before v1; the run is therefore a targeted diagnostic with promotion disabled. The builder reconstructs exactly two \texttt{cwe\_ids} rows and one \texttt{references} row, removes prior labels and rationales from the worklist, and freezes 12 complete official responses. For \texttt{CVE-2024-8020}, a source gate parses the PyTorch Lightning \texttt{2.3.2} \texttt{post\_state} handler and requires the \path{POST /api/v1/state} route, request-body parsing, direct \texttt{body["state"]} indexing, and no local \texttt{try}. For \texttt{CVE-2023-4304}, the gate combines the official CWE-840 mapping restriction with the changed executable lines in the fixing patch and refuses to infer authorization semantics from name/email validation. For \texttt{CVE-2023-32187}, two predeclared reference-resource profiles retain exact successful HTTP resources or narrowly repair only the malformed suffix after a complete CVE Bugzilla identifier. The independent verifier reconstructs the cohort, evidence hashes, source AST, patch markers, and both reference-set relations without importing the analyzer.

We finally perform a post-hoc, read-only staged-frontier audit over the immutable A-E request logs and sealed merge summaries. Request identity is the ordered sample-ID payload; excess request attempts are counted per exact payload, while split retries may establish row completion without repairing attempt-level provenance. The audit separately counts request row-attempts, successful reviewer-row decisions, and token usage recorded on successful responses. It reports marginal candidate yield for A/B strict consensus, reviewer C, and evidence reviewers D/E against the unchanged 0.982 combined-coverage target. A predeclared stop rule rejects further same-model escalation only when both prior advancement gates fail, combined coverage remains below target, and the targeted residual audit promotes no row. This operational audit cannot create labels, impute missing error reasons or token usage, reduce the real-person review queue, or estimate accuracy.

After the D/E evidence stage leaves all 28 affected-version rows unresolved, we run a separate structural work-allocation audit. It reads the sealed blind rows and the previously frozen official-edge result, computes project-family and edge features, ranks repeated families, and only then attaches D/E labels as diagnostics. The fixed rules produce 16 families and select Mattermost ahead of LF Edge EVE and Hutool without using reviewer labels. The two fixed Mattermost rows share stable product/module spans but add Go pseudo-version and legacy-coordinate claims. Release-graph v1 stops before set analysis when the GitHub Releases API reaches its pagination cap; v2 also stops before manifests because two supplied boundary versions are not GitHub Release objects in the frozen prefix. An explicit post-failure v3 therefore tests a narrower 19-token official Git-tag manifest domain. Every tag must bind the current \path{server/v8} module; each pseudo-version must bind an exact commit, timestamp, module manifest, and determinate tag-to-commit ancestry for all 19 tokens; the legacy coordinate additionally requires a total exact-tag manifest mapping. Diverged histories and HTTP 404 records remain unresolved. Family advancement is fixed at 2/2, and a separate cache-only verifier reconstructs the full source inventory and every gate.

We next evaluate the second-ranked LF Edge EVE family under a separately frozen v1 mechanism contract. Protocol discovery before v1 is explicitly disclosed: it identifies a 207-token official Git-tag domain from \texttt{3.0.0} through \texttt{10.1.0}, the nested \path{pkg/pillar} Go module, the non-Go \path{pkg/vtpm} build component, and the paths changed by the two pseudo commits. This makes the run ineligible for candidate promotion or success-rate interpretation. To avoid hundreds of API comparisons, the analyzer freezes the official refs and a content-addressed Git pack containing every commit reachable from the 207 tags and pseudo commits plus the required tag, tree, manifest, and build objects. The cache-only verifier imports the pack into a new bare repository, validates object hashes, rebinds all tags to the official ref snapshot, and recomputes ancestry. Each row additionally requires an actual root Go module for the structured coordinate, advisory-to-repository and component binding, pseudo-path coherence, total main/LTS ancestry, and patched-anchor ancestry. Diverged histories, repository-name similarity, current UI prose, and timestamps cannot fill a missing edge. The family threshold remains 2/2 and all outputs are non-human diagnostics.

The third-ranked Hutool family is evaluated under a snapshot-extensional Maven v1 contract. Its disclosed pre-v1 discovery observes equal Maven Central catalogs for the aggregate \texttt{cn.hutool:hutool-all} artifact and the two advisory components, plus the aggregate structure at three critical versions; candidate promotion is therefore disabled regardless of the result. The stable domain is the exact intersection of the three cached catalogs after retaining only numeric \texttt{MAJOR.MINOR.PATCH} tokens and excluding five fixed milestone tokens. Every catalog must bind its requested coordinate and expose the same stable domain. Source \texttt{hutool-all} POMs at \texttt{5.8.19}, \texttt{5.8.21}, and \texttt{5.8.22} must declare same-parent-version dependencies on both components and a shade goal; the corresponding Maven Central aggregate JARs must contain compiled core and JSON classes. Catalog equality establishes total token correspondence over the frozen domain, while POM/JAR containment is claimed only at the three inspected anchors. Introduced-at-zero GHSA ranges are then evaluated over the cached component catalogs and unioned before comparison with the NVD product set. The family threshold remains 2/2, an independent verifier recomputes all sets from the hash-bound cache without network access, and a passing mechanism requires a later blind cohort rather than promotion on these revealed rows.

We subsequently reuse this frozen mechanism on a CVE-exposure-disjoint Hutool cohort from the same aligned snapshot. The exclusion parser projects only CVE identifiers from prior artifacts; baseline, reviewer, consensus, and candidate fields do not enter selection. The cohort builder scans all 8,066 matched rows, excludes the union of 1,967 previously exposed CVEs, and retains six Hutool rows that satisfy one of two predeclared routes: direct product-to-aggregate correspondence or product-to-aggregate-component correspondence. It seals the six rows before candidate computation. The mechanism predates this availability audit, but the cohort's availability and structural shape are observed in the same snapshot before sealing; we therefore preclude promotion and treat the run as a retrospective external application, not a future time holdout or human evaluation. The verifier independently reconstructs the exposure union, cohort, release sets, relations, and no-promotion gate from the authoritative aligned input and hash-bound v1 cache.

For the references candidate, we audit all 56 rows changed from representation discrepancy to incomplete. The builder recomputes every derived status, binds the exact changed-CVE set, and freezes 118 URL probes across 57 identity groups. All 56 rows enter both reviews; a 28-row automatic network-positive subset is not exempted. The reviewer worklist exposes neutral group IDs, original NVD/GHSA URLs, and frozen probes, but not the transformation, candidate identity, trigger stage, subset direction, automatic certificate, or expected label. Reviewer E inspects probes before URL structure, while reviewer F reverses that order. A protocol pilot is retained but excluded because its validator enforced a confidence rule absent from the sealed prompt. Revision 2 seals both complete role prompts and uses fresh ephemeral runs. Strict consensus requires matching determinate verdicts and group decisions with no additional-review request. The resulting narrower profile is selected after this audit and remains post-hoc, same-model-family, and non-human.

A separate blank human-review packet binds the same revision-2 seal, raw URLs, frozen probes, group order, and verdict-to-status contract for all 56 rows. It schedules a primary annotator, a different independent reviewer, and an author resolution/signoff. Because the non-human disagreement reflects two resource-identity constructs, each stage must explicitly select an underlying-content, frozen-HTTP, or written custom identity definition before recording group decisions. The 24 encoded-line rows are prioritized but not prelabelled. At the current snapshot all 56 rows are pending and none is signed; the packet and readiness artifact remain \texttt{label\_is\_human=false} and cannot support a human-gold claim.

For a non-human CWE candidate diagnostic, we additionally audit the complete 17-row impact set of \texttt{taxonomy\_v1}. The blind worklist and 34 current/candidate predictions are sealed before either reviewer file exists. Two fresh Codex runs receive only the blind CWE 4.20 context in isolated workspaces. Sixteen affected CVEs are disjoint from the 300-row primary seed. Strict consensus requires agreement on set relation, discrepancy label, and taxonomy-support verdict, with neither reviewer requesting additional review. This design tests the complete rule impact set but remains impact-selected, same-model-family, and ineligible for a human-gold claim.

After the first audit, its six unresolved rows and three strict candidate regressions form a post-hoc priority subset for an evidence-enhanced secondary audit. Before two additional reviewer files exist, a new builder freezes up to five ranked references listed by NVD or GHSA per row, the secondary prompt, the first-stage artifacts, and code hashes. The secondary blind worklist contains neither method prediction nor prior reviewer label. One new run reads frozen evidence before the CWE mapping; the other identifies the concrete mechanism first. A determinate decision must cite a literal substring from a successfully fetched snapshot. Secondary strict consensus requires agreement on relation, label, taxonomy support, and concrete-mechanism verdict. This stage diagnoses whether additional evidence resolves the selected boundary cases; because selection follows the first audit and evidence availability is uneven, it is not an independent confirmation.

For the severity adjudication prototype, the analysis uses the 80-sample severity seed and its evidence cache. The label distribution in the silver-v2 output is 75 factual\_conflict, 3 representation\_discrepancy, and 2 temporal\_discrepancy. The adjudicated-source distribution is 49 both, 26 NVD, 3 GHSA, and 2 abstain; confidence is 69 high and 11 medium.

For affected\_versions, the analysis uses the 100-sample manual-check set and the evidence-aware silver-v2 output. The silver label distribution is 47 factual\_conflict, 28 incomplete, 17 representation\_discrepancy, 4 equivalent, and 4 uncertain. The adjudicated-source distribution is 45 NVD, 35 both, 11 GHSA, 7 abstain, and 2 neither; confidence is 34 high, 64 medium, and 2 low.

We additionally froze a second 100-row affected\_versions holdout from the current 651 deterministic factual-conflict candidates. Selection used a fixed SHA-256 rank after excluding every CVE in the earlier 100-row development cohort, leaving 551 eligible CVEs. Evidence was collected in a dedicated cache; all 100 rows had at least one nonempty fetched-text record. A blind worklist exposed source values, context, and evidence but excluded previous annotations, predictions, and correctness fields. Predictions from 18 fixed methods were sealed before either reviewer decision file existed. Two Codex reviewers independently assigned both a discrepancy type and a supported source from the same blind worklist. The strict merge retained only determinate rows with exact joint agreement and valid URL-level support. These labels remain \texttt{label\_is\_human=false}: the design is development-disjoint and leakage-controlled, but it is neither human gold nor independent-human validation.

The preregistered v1 endpoint was source-selection accuracy over all strict joint-consensus rows. After unsealing, we found that this endpoint mixed factual-conflict source adjudication with non-conflict discrepancy typing. We therefore report the original endpoint unchanged and a clearly marked post-hoc split between factual-conflict and non-conflict rows.

We then froze a v2 cohort after excluding the 200 CVEs in both earlier affected\_versions cohorts, leaving 451 eligible deterministic conflict candidates before fixed-hash selection of 100 rows. Its dedicated snapshot contains 564 evidence records and usable fetched text for every row. Before either new reviewer file existed, we sealed 300 type predictions from three methods and 1,900 source predictions from 19 methods, together with method and protocol code hashes. V2 preregisters two endpoints: full discrepancy-type accuracy on strict type consensus, with method abstention counted as incorrect, and source accuracy only on rows with both strict factual-conflict type and strict source consensus. The reviewers cite literal frozen snippets through structured evidence claims. Every v2 decision and consensus row remains \texttt{label\_is\_human=false}.

After v2 was unsealed, we performed a clearly marked development-only leave-one-cohort-out diagnostic over the Phase D, v1, and v2 non-human labels. Gold-blind structured features exclude CVE identifiers, sample identifiers, discrepancy labels, and adjudicated sources. Balanced logistic regression and a shallow decision tree are trained on two cohorts and tested on the third. The advance gate requires a candidate to produce strictly more correct predictions than the best named comparator in every held-out cohort. Because the feature and model families were selected post-unseal, this diagnostic can reject a candidate but cannot confirm a method.

The severity diagnostic baselines are \texttt{prefer\_nvd}, \texttt{prefer\_ghsa}, \texttt{latest\_published}, and \texttt{evidence\_score\_baseline}. The affected\_versions silver-label diagnostic comparison uses the same fixed-source and latest-published baselines plus \path{version_token_support_baseline}, a simple fetched-text matcher that predicts support when evidence snippets mention at least one affected-version token from a side. The older URL-only LLM draft is not reused as an adjudication label set, and silver-v2 is evidence-aware but still not gold.

To prepare a human audit path for RQ3, the repository includes blank templates under \path{data/annotations/rq3/gold_audit/}: 80 severity rows and 100 affected\_versions development rows. These templates preserve the source evidence, source identifiers, evidence summaries, and silver-v2 annotations as provenance, but all \texttt{human\_audit} fields are still draft/blank. The guarded evaluator \path{experiments/rq3_adjudication/evaluate_rq3_human_audit.py} refuses the current templates and writes no gold-audit metrics until valid final human rows are present. Both CVE-disjoint holdout consensus artifacts also require real human signoff before they can support a human-gold claim.

The current paper artifact summary is written to \path{results/paper_cose/cose_artifact_tables.md} and \path{results/paper_cose/cose_artifact_tables.json}. Those files aggregate the reproducible counts used in the paper.

\input{sections/05_results.tex}
\section{Discussion}

The current artifacts support three practical observations. First, post-alignment discrepancy is not uniform: severity and affected\_versions are where deterministic factual-conflict candidates remain most visible, while published is often a temporal or representation issue, references is frequently incomplete, and cwe\_ids is mostly equivalent with a smaller incomplete/conflict tail. Second, the RQ2 diagnostic shows why a binary same/different view is too coarse: most binary-different field instances are assigned to non-conflict categories by the deterministic baseline. Third, the diagnostic evidence-source profile is uneven by source class and field: severity has broad fetched-text availability, while affected\_versions diagnostic samples include a larger number of GitHub and vendor evidence URLs with lower fetch success.

The key implication is operational rather than purely taxonomic. A downstream consumer should not react to every binary mismatch with source selection. Equivalent and representation differences mainly call for normalization, incomplete records call for augmentation or coverage analysis, temporal differences call for timeline interpretation, and only residual factual conflicts call for evidence-constrained adjudication. This routing view is the main method claim that can be defended before gold labels exist: post-alignment fusion needs a discrepancy type before it needs a winner.

The CWE full-impact audit provides the strongest current positive RQ2 candidate signal. After a post-hoc evidence-enhanced pass over the original priority rows, its direction favors treating fully related ancestor/descendant assignments as representation differences on 10 of 14 primary-seed-disjoint strict rows. Additional evidence resolves four of the original six undecided rows, including two OS-command cases, but it also expands the concrete boundary to four factual conflicts despite complete taxonomy connectivity. Two mechanism-specific rows remain unresolved. Evidence therefore improves candidate coverage from 11/17 to 15/17 without making the rule automatic: taxonomy ancestry still cannot establish that the specific CWE applies to the CVE. The post-hoc selection, uneven snapshot availability, same-model reviewers, and absent human signoff make this a prioritized human-audit result rather than a production-switch result.

The references audit exposes a complementary construct boundary. Exact redirects and preserved advisory or bounty identifiers yield 32/32 strict agreements across four rule families, but every encoded-line row splits the reviewers. Both interpretations are coherent: an encoded \texttt{\%23L} suffix can be intended as a line anchor on one file, yet it is also literal path text in the observed request and returns 404. The narrower audited profile avoids choosing the interpretation that best improves the candidate metric. Its exact 32-row impact alignment is useful development evidence, but only a human-defined resource-identity contract can decide whether malformed line references should be normalized, repaired, or retained as distinct failed references.

The three-row residual audit shows the same boundary at case scale. Source code can make a local weakness mechanism concrete, as the unhandled PyTorch Lightning request-body access does for CWE-248, without turning that development judgment into gold. Official taxonomy metadata can also invalidate an abstraction level for vulnerability mapping, as with CWE-840, while still failing to identify the correct concrete weakness from a narrow patch. Finally, the K3s malformed references produce different labels under two coherent resource identities. More evidence is therefore not interchangeable with an approved construct: code can support a mechanism, but taxonomy and URL identity still require an explicit review contract.

The staged frontier adds an operational stopping argument. A third same-model pass resolves 64\% of its selected rows, but two evidence-enriched passes resolve only 16\% of the next selected remainder and 8\% per reviewer-row decision. These selected-stage ratios are not comparable accuracy estimates, yet they show that generic same-family escalation has reached diminishing candidate yield under the current contract. Because the fixed target is still missed and later targeted diagnostics cannot be promoted, the defensible next investment is construct approval plus independent real-person review, or a later frozen snapshot, rather than another vote on the revealed rows.

The same artifacts also show why the paper must stay cautious. RQ2 now has a large non-human typing-stability diagnostic, but still not a gold-label agreement study: its two pass roles share one model/config, comprise many ephemeral sessions, and cannot compare candidate profiles because all predictions are identical. In this non-human diagnostic, the post-unseal analysis shows that the severity result largely compares a label-only baseline with a richer score/vector/version review construct, while affected\_versions exposes 25 concrete unbounded-record projection losses. The resulting 1,121/1,147 post-hoc fit is explicitly in-sample, and the cross-protocol gate then rejects it: severity moves by -23 and -7 rows on the older primary/review slices but +165 on fresh strict labels, while affected inputs are not preserved consistently. This reversal makes protocol drift and input-projection sensitivity a stronger finding than candidate accuracy. Severity remains non-human, the complete 1,250-row human packet is blank, and the corpus has no human-approved field contract yet.

The later isolated snapshot removes the earlier profile-identifiability failure but not the validation boundary. Three CWE rows differ between current and the taxonomy candidate, yet only one receives strict dual-Codex consensus in the sealed main review. Frozen source evidence then makes both fresh reviewers agree on the candidate's granularity-only direction for all three cases, including the two previously non-strict rows. This strengthens the case-level mechanism explanation, not the profile metric: the rows are selected after the differences are known, and the reviewers still share one model family. The two rejected protocol attempts also show why exact output interfaces belong to the evidence contract: malformed confidence or path representations must fail closed rather than be silently normalized into agreement. More importantly, the strict event-time universe is empty: all 250 reviewed rows are snapshot-external rather than event-time-external. The result therefore supplies three prioritized CWE cases and a reproducible availability finding, not a generalization or method-gain estimate. It also illustrates why a later download is insufficient by itself; temporal validation must bind source event timestamps, frozen profiles, and untouched records simultaneously.

The 50-row CWE audit provides a stronger field-control check but reaches the same claim boundary. Hiding the three differences among 47 controls preserves their candidate direction and yields 48/49 candidate agreement versus 45/49 current agreement, while exposing one evidence disagreement and one profiles-equal miss. This reduces the most direct case-selection cue, but it does not undo that CWE was chosen after profile unsealing and that all reviewers use the same model family. Its failed protocol versions are also substantive findings. Literal set inclusion does not by itself imply semantic incompleteness when the extra CWE may describe a different mechanism, and a citation is not auditable unless its quoted text is a literal frozen substring before the decision is accepted. The V3 pre-write checks make those assumptions executable. They improve auditability, not external correctness, so the sealed +1/231 comparison remains the only main profile metric.

The final 16-row evidence stage reinforces the stopping decision across fields. Every selected row has a successfully retrieved non-empty source record, yet only both references and two affected\_versions rows become strict; severity remains 0/2 and affected\_versions 2/12. The bottleneck is therefore not HTTP availability alone. It is whether the supplied records establish the intended CVSS comparison, URL resource identity, or complete product/package/range relation under a shared contract. Reaching 238/250 staged coverage while failing the 0.40 selected-row resolution gate is not a near-pass to be rescued by threshold relaxation. It is evidence that another correlated vote over the same revealed cohort is unlikely to resolve the remaining construct choices.

The label-free envelope sharpens the practical interpretation further. Because 98.8\% of the sealed predictions are identical, no human labeling of this cohort can produce more than a 1.2-point absolute accuracy difference between current and the CWE profile. The six profiles reduce to two prediction vectors, and every cross-vector pair differs on the same three rows. Even the most favorable possible gold assignment gives an exact two-sided McNemar p-value of 0.25, so no labeling of this cohort can support rejection at alpha 0.05. Six same-direction correctness-discordant rows are only the theoretical minimum for any rejection; exact 0.80 power requires 49, 20, or 12 effective discordant rows under candidate-win probabilities 0.70, 0.80, or 0.90. This label-free diagnostic is not evidence that either profile is accurate, and those power values are conditional assumptions rather than a sample-size commitment. The three differential rows determine the paired sign but cannot validate the taxonomy over the other 247 rows or establish absolute performance and construct error. A future confirmatory cohort should therefore require event-time eligibility, enough naturally occurring sealed profile differences, and an a priori power calculation tied to a human-approved outcome contract.

The full eligible-universe census changes the planning interpretation of that calculation. The six profiles are all distinct over the full revealed universe, but only 34 of 29,740 field instances differ: 29 CWE rows and five references rows under the broader references profile, reduced to three under the audited profile. Thus the sample's 3/250 rate should not be extrapolated as stationary; it comes from a current-status-stratified allocation, and the full prediction changes are concentrated in rare rule-trigger strata. A future strict event-time evaluation should first census sealed predictions over its eligible universe. It can then freeze a disagreement-enriched review layer, including all differences when feasible, for paired method comparison, while retaining a separately probability-sampled control layer for absolute accuracy and prevalence. CVE-level grouping and at most one primary comparison row per CVE should prevent pseudo-replication. Gold labels and reviewer identities must remain hidden from prediction census and selection. On the current revealed snapshot, references profiles have only five or three possible current-versus-candidate discordant CVEs, below the six-row extremal significance minimum, whereas CWE and combined profiles have 29--34 prediction differences but unknown correctness discordance. This is design information, not power or method evidence.

The second label-free acquisition makes the temporal bottleneck observable rather than treating the empty strict tier as a generic waiting condition. NVD adds 39 records published after the profile freeze, while the reviewed GHSA snapshot adds none; therefore no bilateral post-freeze match exists and the frozen field views remain unchanged. This supports a source-aware readiness gate: collection time alone cannot start confirmatory review, and unilateral source growth should not be counted toward a bilateral event-time denominator. The correct action is to retain the sealed profiles, rerun the same label-free availability audit after GHSA updates, and freeze a cohort only when the predeclared minimum is met. This monitoring result is not a temporal validation result because it contains no eligible row and no correctness label.

Auditing the complete revealed difference set sharpens, but does not remove, that boundary. The 25-to-1 CWE direction shows that the taxonomy profile's changed rows are usually compatible with the same-model evidence contract, while the one current-direction row and three unresolved rows show that ancestry is not sufficient as an automatic rule. Conversely, the references result shows why field status alone is coarser than resource identity: reviewers can agree that a target GHSA alias is shared yet disagree on whether NVD and third-party pages belong to that same underlying document group. Requiring complete partition identity exposes this disagreement rather than granting a candidate win from one selected pair. Future evaluation should therefore freeze the resource ontology and adjudication unit with real reviewers before profile comparison, not derive them from whichever partition improves the metric.

The third blind pass does not remove this limitation. Its fixed majority rule resolves 66 of the 103 non-human A/B non-strict rows, but misses both advancement thresholds and leaves 37 non-human candidate rows, including 28 affected\_versions cases. Within these non-human outputs, triple uncertainty on 17 rows and three-way qualified-label splits on seven show two distinct bottlenecks: some inputs lack enough evidence for a determinate decision, while others expose taxonomy or field-contract ambiguity even when the model is confident. A fourth same-family vote would add correlated evidence rather than resolve either construct. The next useful step is evidence enhancement and real-person semantic adjudication, not threshold relaxation.

The evidence-secondary result distinguishes generic evidence availability from construct-complete evidence. Frozen advisory, commit, issue, and release-page text is available for 28 of the 37 non-human rows and supports six strict D/E candidates outside affected\_versions. It resolves no non-human affected\_versions row: both reviewers still abstain together on many cases because text fragments do not establish complete product/package identity, all range endpoints, branch/backport scope, or a total component-to-product release mapping. The shared 19 non-human uncertain decisions make exact agreement a misleading success measure. For this field, the next evidence unit must be a project-specific release, manifest, lockfile, or vendor mapping graph under a human-approved range contract; adding more generic URLs or correlated votes does not address the missing edge.

The disjoint non-human development calibrations make the construct boundary more specific. They support a candidate distinction between same-CVSS-version base-metric conflict and cross-version representation, and they repeat the missing-score and one-sided-unbounded rules on disjoint rows. They also show that high aggregate agreement is insufficient: v2 has perfect exact agreement, yet its only non-strict row is a shared abstention. Official component evidence alone does not resolve that abstention, and the current-coordinate-only diagnostic abstains under its frozen contract. The historical non-human graph shows what additional structure is needed: product-build dependencies map product versions to legacy artifact versions, source identity binds the migration edge, and only then can prerelease sets be ordered. That graph yields a post-hoc non-human \texttt{incomplete} candidate for the XWiki case. A generic fail-closed schema then binds all eight structurally selected equal-boundary cases across three ecosystems, showing that package identity, component membership, and an explicit artifact alias can be represented by one evidence contract. The non-equal audit reveals that released-version set equality is not the same construct as syntactic interval containment. Freezing the assistant's extensional choice and applying it to two InLong components makes that alternative executable, but the representation candidate still conflicts with both reviewers. The full-aligned-data stress test then shows a separate boundary: heterogeneous package claims cannot be unioned merely because registries enumerate both sets. Parallel distributions, dependency constraints, and coordinated components require different evidence edges, and none of the initial NuGet, PyPI, or crates.io rows has a total product-release map. The post-no-go Deno result narrows this claim: a dependency range is insufficient, but an exact committed build lock can recover the product edge and expose five runtime-only affected releases. This is evidence for an edge-specific mapping rule, not for treating all lockfiles as authoritative or for resolving the field's extensional/intensional choice. Real people must still choose the range semantics and approve acceptable lockfile, build-manifest, or vendor mappings for each edge class.

The Mattermost family audit shows that project-specific evidence is necessary but still not sufficient. Exact tag manifests establish the current module identity for every fixed token, and the pseudo-version timestamps identify unique commits, yet most product tags are on incomparable backport branches. Treating timestamp order as release ancestry would manufacture a total mapping that the commit graph does not support. The historical module path adds a second unresolved union edge. The defensible outcome is therefore 0/2 abstention, not a representation or incompleteness label. Together with Deno, this isolates two different edge classes: exact build locks can recover one product-to-component sequence, while repository identity plus a pseudo commit does not recover a multi-branch release projection.

The EVE family makes the failure more specific. A GHSA package string can name a repository that contains several nested Go modules and non-Go build components without being a module itself. In addition, a syntactically valid pseudo version identifies a commit boundary but does not guarantee that the commit modifies the vulnerable component: one EVE pseudo changes a filesystem setup script outside \path{pkg/vtpm}. The public advisory API and repository advisory UI also expose different levels of version detail. A content-addressed commit graph can remove API-rate and reproducibility barriers, but it cannot turn those interface, component-identity, or branch-divergence gaps into a total product set. The 0/2 result therefore supports field-specific abstention and evidence-interface preservation, not a claim that the API, UI, NVD, or GHSA is factually wrong.

Hutool isolates a different condition: a registry can expose a coordinated aggregate/component release line without supplying a human-approved meaning for an unbounded affected range. Equal Maven catalogs make the release tokens comparable, and the anchor POM/JAR evidence prevents a name-only product-to-component projection. Under the declared snapshot-extensional semantics, both NVD sets are strict subsets of the GHSA component unions. That mechanism result is useful, but it does not establish that the current 209-release registry snapshot is the intended historical universe or that every aggregate JAR was inspected. Since the mapping shape was known before v1, the defensible consequence is a new-cohort requirement, not promotion of the two revealed \texttt{incomplete} candidates.

Applying the frozen mechanism to six CVE-exposure-disjoint Hutool rows separates mechanism reuse from repeated judgment on the original two cases. The six projections include both strict-subset and equal-set outcomes, so the contract does not mechanically assign every eligible row to \texttt{incomplete}. This strengthens the evidence that the cached token correspondence is executable across more than the development pair. It still does not establish temporal generalization or accuracy: cohort availability was discovered in the same aligned snapshot, labels remain Codex development candidates, and promotion is disabled.

The affected\_versions field remains a special case because version-system mismatches can still produce diagnostic conflict signals. OSV/GHSA-style affected records separate package identity and range events, and recent affected-version work shows that this field remains difficult even when tools use multiple forms of advisory evidence \cite{OSVSchema2026,GHSARepo2026,Chen2025AffectedVersions}. The current baseline has been tightened and NVD \texttt{vulnerable=false} CPE matches have been filtered, but the residual 651 factual-conflict candidates still need manual spot checks before they can be treated as stable field labels. The repair changed only 10 full-corpus classifications and no frozen-cohort decision, but that stability is an observed input-impact result, not evidence that package identity or range semantics are solved.

The development evidence-aware pass for affected\_versions supports an evidence-availability statement, a silver-label distribution, and a silver-only baseline comparison. The v1 CVE-disjoint holdout adds a stronger leakage-control check: its cohort, evidence cache, blind reviewer input, and predictions were separately frozen, and the predictions predate both decision files. This improves auditability, but it does not convert Codex decisions into human gold. Exact joint consensus covers only 35/100 rows, and the reviewer agreement values indicate substantial unresolved ambiguity.

More importantly, the v1 post-unseal diagnostic changes how the affected\_versions result should be interpreted. The preregistered all-strict diagnostic score rewards a method both for recognizing non-conflict cases and for selecting a source within true factual conflicts. In this diagnostic, token methods score mainly through the former because representation discrepancies tend to map to \texttt{both}; on the 16 strict factual conflicts, no proposed method exceeds fixed GHSA preference or recency. V1 therefore exposes the need to separate discrepancy typing from conditional source adjudication.

V2 preregisters that separation on a new cohort, so its negative result is more informative. In this non-human diagnostic, the type candidate is correct on its three determinate predictions but abstains on 38 of 41 strict type rows; counting abstention as incorrect makes its full accuracy substantially lower than both full-coverage comparators. On the nine non-human strict-source rows, the branch source head reaches 2/9, while fixed NVD preference reaches 6/9. Thus endpoint separation fixes the interpretation problem but does not rescue method performance. The defensible finding is a protocol correction plus a failed, highly selective method candidate, not an affected\_versions adjudicator.

The post-v2 diagnostics sharpen that conclusion rather than overturning it. A more permissive task-separated type candidate and an authority-filtered source candidate do not improve consistently across Phase D, v1, and v2. A leave-one-cohort-out classifier improves the pooled type count but regresses below the legacy comparator on held-out Phase D, while its source endpoint is below the branch baseline both pooled and on Phase D. This cohort sensitivity is why we do not spend another sealed cohort on the current candidate family. The next methodological requirement is a human-audited source-authority contract and version semantics that first demonstrate cross-cohort stability on development data.

The evidence audit also separates reviewer independence from evidence independence. Five of nine v2 strict-source rows use the same single cited URL for both reviewers, and only two have primary/ecosystem evidence cited by both. Agreement under a shared evidence snapshot can therefore reflect a common evidence bottleneck. Future source labels must record source ownership and distinguish upstream maintainer, ecosystem database, disclosure, database record, and secondary evidence before a strict source decision is accepted.

The sensitivity diagnostic reinforces that caution. Severity predictions change for fewer samples across evidence-score thresholds, while affected\_versions predictions change for most samples when the minimum version-token support threshold is raised. That instability is not a failure by itself, but it shows that token presence is a brittle proxy for source support unless package identity and version-range semantics are modeled explicitly.

The case sketches reinforce the same boundary. The severity examples show both the value and the weakness of a simple evidence-score rule: exact score/vector matches can recover source support, but they can also miss lower-salience evidence for the other side. The affected\_versions examples show why this field needs extra care around package identity, patch-level granularity, and ecosystem-specific version schemes before a conflict detector can be treated as reliable.

The silver error-mode diagnostic gives a concrete audit target for the next iteration. Severity mismatches are often source-split differences where the silver label marks both-source support but the baseline marks only one source, or the reverse. Affected\_versions mismatches are more structurally diverse, and the current 100-sample diagnostic highlights many no-overlap package-name cases. This does not validate the silver labels as gold, but it does identify where a future human audit and semantic version-range adjudicator should focus first.

This positioning also constrains the related-work claim. The contribution should not be framed as superseding vulnerability text discrepancy mining, general truth discovery, or affected-version analysis. Its defensible space is narrower: a post-CVE-alignment workflow that makes structured-field disagreement auditable and routes only the residual conflict cases to source-support analysis with abstention.

The main discussion point should therefore be selective and modest: the current evidence-aware results are diagnostics within non-human label setups. The strongest supported contribution is the auditable type-first pipeline, including explicit abstention, sealed prediction artifacts, the identified task-mixing failure mode, evidence-dependence audit, and a follow-up protocol that evaluates the two tasks separately. Any broader claim about affected\_versions method effectiveness or general vulnerability-field fusion must wait for human gold. A further untouched cohort is warranted only after a revised candidate first passes the documented cross-cohort development gate.

\section{Threats to Validity}

Construct validity is limited by the current label source. The artifact set has deterministic baseline outputs for field typing, silver-v2 development annotations, and strict overlays from blind Codex reviewers, but no human gold. The five discrepancy categories are useful for analysis, yet their boundaries can be subjective, especially between representation\_discrepancy and incomplete for affected\_versions. The older URL-only LLM draft is not treated as an evaluation label set, and the dual-Codex overlays are explicitly marked \texttt{label\_is\_human=false}.

The routing contract itself is also a construct-validity assumption. The method maps equivalent and representation differences to normalization, incompleteness to augmentation, temporal differences to timeline interpretation, and factual conflicts to evidence-constrained adjudication. That mapping is methodologically useful, but it is not yet validated against practitioner decisions or downstream scanner behavior. Until such validation exists, the routing contract should be read as an auditable baseline design, not as a proven operational policy.

Internal validity is limited by evidence quality. The evidence cache stores fetched text snippets and metadata, not a full reconstruction of dynamic pages. Some URLs return 403, 404, or timeout, so the evidence set is incomplete by construction. The evidence-source diagnostic shows that this incompleteness is source-class dependent: affected\_versions has complete fetched text for NVD pages in the sample, but substantially lower fetched-text rates for GitHub commit/repository pages and GitHub advisories. These failures can bias adjudication toward sources that are easier to fetch. The silver-v2 labels also depend on the fetched evidence context; if important vendor or commit evidence is missing, both the silver label and the baseline comparison can be biased.

Reviewer independence is not evidence independence. In the nine v2 strict FC-source rows, both reviewers cite the same single URL in five cases, and only two cases contain primary/ecosystem evidence cited by both reviewers. Three rows collectively rely only on NVD record evidence. The strict overlay may therefore favor evidence that is shared or easier to fetch, even when the reviewer processes are isolated.

The affected\_versions holdouts reduce prediction-to-label leakage but do not remove all dependence. V1 excludes the development CVEs; v2 excludes both earlier 100-CVE cohorts. Reviewer inputs omit prior labels and sealed predictions, yet all rows are sampled from the same deterministic factual-conflict candidate miner. Each reviewer pair also shares the same Codex model family, prompt contract, and evidence snapshot, so their errors are not independent. In these non-human diagnostics, v1 label/source kappa is 0.268/0.392 with 0.35 strict joint coverage. V2 label/artifact kappa is 0.535/0.669, but strict type coverage is 0.41 and strict FC-source coverage over the full cohort is only 0.09. Reported performance is therefore conditional on selective, easier-to-agree subsets.

External validity is limited by coverage. The adjudication diagnostics cover severity with 80 silver samples, affected\_versions with a 100-row development sample, and affected\_versions with two additional 100-row CVE-disjoint holdouts. Both holdouts remain conditional on the current NVD-GHSA corpus and candidate miner. The affected\_versions predictors are not semantic version-range adjudicators, and affected-version identification itself is known to depend on package, range, and ecosystem semantics \cite{OSVSchema2026,Chen2025AffectedVersions}. Neither field has gold-backed metrics. The current numbers should not be generalized to all fields or vulnerability databases without the pending human audit and cross-field experiments.

Conclusion validity is limited by the selective setup. RQ2 currently reports diagnostic rule-trigger counts rather than performance metrics. Accuracy and macro-F1 are reported against non-human labels, and coverage matters because both adjudicator disagreement and method abstention are selective. The v1 factual-conflict/non-conflict split was discovered after unsealing, so its 16-row factual-conflict analysis is post-hoc. V2 preregisters separate endpoints, but its type method predicts only three strict rows and its source endpoint contains only nine strict rows. V1 and v2 are both now unsealed and must not be used to tune a replacement method.

The RQ2 CWE audit seals predictions before reviews and covers the complete 17-row candidate impact set, but it is not a representative sample of the 8,066 aligned pairs. Sixteen CVEs are disjoint from the primary seed, yet the rule family was selected after earlier AI-candidate error inspection, and all reviewers remain Codex runs from the same model family. The evidence-enhanced stage is additionally selected from first-stage disagreements and regressions; only 28/36 reference snapshots were fetched successfully, so evidence availability can affect adjudication. Combining stages raises strict coverage from 11/17 to 15/17, but two rows remain unresolved and four evidence-supported rows favor the current factual-conflict decision. The seed-disjoint bootstrap interval reaches zero and the exact sign diagnostic is 0.1796. These properties permit a bounded candidate-direction statement, not human-gold accuracy, independent confirmation, or a production switch.

The references audit also covers a complete rule-impact set rather than a representative sample. The two reviewers share one model family and evidence snapshot, and the grouped input necessarily reveals which URLs are being tested together even though the candidate transformation and label are masked. Thirty probes return 404, making the resource-identity definition consequential. The first protocol seal is retained but excluded because its validator contained a confidence constraint absent from the prompt; revision 2 uses fresh reviewers and complete sealed role prompts, but it was created after pilot output existed and is not fully preregistered before any model response. The 32-row audited profile is then selected after revision-2 unsealing. Its candidate-label gains are therefore post-hoc development diagnostics, not independent or human-gold performance.

The references human-review packet removes the hidden-contract problem by requiring each reviewer and resolving author to record the resource-identity definition, but it does not itself resolve construct validity. All 56 rows are still pending and no real human has signed one. Its 24-row priority flag identifies the known definition-sensitive subset without supplying either Codex verdict; this is workflow readiness, not human validation.

A separate current-snapshot RQ2 typing-stability cohort excludes all 717 CVEs previously exposed through development, audits, or holdouts and seals six executable profile predictions for 1,250 globally unique CVEs. After those exclusions, all six prediction columns are identical, so the completed review estimates only baseline stability and cannot identify candidate-profile gain. In this non-human diagnostic, the A/B roles are composed from 28 and 67 disjoint ephemeral sessions using the same model/config, not two persistent or human annotators. A later immutable-log audit identifies three excess B requests and 90 retry row-items. Every requested row eventually has an accepted output, but the old log lacks response-error events and rejected JSON; for one duplicate payload, it cannot identify which attempt failed. The new artifact records this missingness without backfilling an error reason. Accepted rows pass the sealed merge contract, but this retry/provenance gap and same-model dependence limit interpretation. Independent comparison of the references and CWE candidates requires a new time cohort collected after profile freeze, and the current consensus still requires human validation.

Both acquisitions are post-seal by collection time but not by bilateral event time. The second acquisition adds 74 NVD rows, including 39 published after the freeze, but adds no reviewed GHSA row and leaves the one-GHSA matched field views byte-identical. No aligned record therefore has both NVD and GHSA publication timestamps later than the profile seal, so the 250-row fallback remains explicitly snapshot-external and no new cohort is frozen. This source asymmetry means that repeated acquisition can audit readiness but cannot by itself establish temporal validation. The fallback's two passes share the same model, prompt, executable, and evidence projection; disjoint sessions and reverse order reduce state and ordering dependence but not model-family dependence. Only three sampled rows distinguish profiles, and two lack strict consensus in the sealed main review. The observed +1 candidate paired count is consequently sparse, selective same-model development evidence. A subsequent evidence-secondary diagnostic selects exactly those three known differences; two fresh same-family passes reach 3/3 strict candidate-direction agreement, but this post-selection result cannot be substituted into the main 250-row metric. V1 and v2 are excluded wholesale after fail-closed checks reject a confidence/uncertainty mismatch and a nonliteral CWE path, so the accepted v3 output avoids silent repair but remains dependent on an interface refined after failure. None of these stages can establish future-snapshot generalization, a temporal effect, human correctness, or a production-switch benefit.

The all-50 CWE evidence audit removes explicit profile-difference flags from reviewer inputs and covers the entire sealed field, but it remains post-hoc because field selection and all three contract versions follow A/B and profile unsealing. Its 134/135 successful URL snapshots are broad but not proof that the retrieved pages contain complete mapping evidence. V1 confounds literal subset with semantic compatibility, and V2 discovers the literal-quotation gap only at merge; both are archived and excluded, while V3 is accepted only after the interface is refined. The 20 sessions are mutually disjoint but still use one model family. Consequently, 49/50 strict coverage and the 45/49 versus 48/49 comparison are useful construct and control diagnostics, not a preregistered accuracy estimate. They cannot replace the sealed 185/231 versus 186/231 result, establish human correctness, or justify production promotion.

The final 16-row non-CWE stage is also outcome-selected after A/B unsealing. Its selector and thresholds are fixed before its own evidence retrieval, its 50 frozen records all fetch successfully, and G/H sessions are disjoint from prior stages, but G/H still share the same model family and prompt lineage. The staged 238/250 candidate preferentially retains already strict rows and adds decisions only from difficult revealed subsets; its 0.952 coverage and 188/238 versus 191/238 profile agreement therefore cannot estimate full-cohort accuracy. The latter difference is inherited from the earlier CWE stage, since the four new non-CWE decisions do not distinguish profiles. Successful retrieval also does not prove completeness of package mappings, range endpoints, CVSS semantics, or resource identity. The fixed resolution gate fails at 4/16, so neither the 12 remaining abstentions nor the no-go may be converted into human truth.

The paired-outcome and paired-test diagnostics are label-free and exhaustive over the frozen five-label universe, but their 125 assignments are not equally likely empirical outcomes. Treating the 40 candidate-better, 40 current-better, and 45 tied cases as probabilities would impose an unsupported uniform prior. The +/-1.2-point bound applies only to full-cohort accuracy. The exact-test result proves only that no labeling of these three prediction differences can reject the paired null at alpha 0.05; it does not establish method equivalence. The future-cohort counts additionally assume independent sampling and a stationary 3/250 prediction-difference rate, while the power table conditions on an assumed candidate-direction probability among correctness-discordant rows. Neither diagnostic bounds macro-F1, category-specific recall, construct validity, the effect of a revised profile, or the true power of a future cohort without those assumptions.

The eligible-universe census removes the unsupported stationary 3/250 extrapolation for this revealed snapshot, but introduces no correctness evidence. Its 34 union differences are deterministic outputs from profiles developed with earlier diagnostics; the universe is snapshot-external rather than strict event-time, and its contents and rule-trigger strata are now revealed. The census cannot be converted into a confirmatory cohort or an unbiased estimate for a later snapshot. As a label-free diagnostic, current-status stratification explains why the selected slice is not self-weighting, but does not by itself explain every within-stratum deviation in a small sample. Most importantly, this diagnostic's 29 CWE prediction differences and five or three references differences are not correctness-discordant counts. A disagreement-enriched future review cannot support absolute performance or prevalence without a design declared before labels, CVE-level grouping, and a separately probability-sampled control layer.

The two complete-difference audits add labels only after the universe and outcomes are revealed. Their completeness prevents within-impact sampling loss but does not make the selected strata representative, event-time eligible, or independent of profile development. The CWE reviewers share a model family and all rows expose the same disjoint taxonomy structure; the small conditional p-value is therefore not a preregistered significance result. The references worklists hide direct source fields but raw URLs can reveal CVEs and the alias class, and complete-partition consensus is sensitive to how related advisory, NVD, and third-party pages are ontologically grouped. Revision 1 also contains a disclosed merge-code failure; revision 2 is valid only as a repaired development protocol created after that failure. Neither result is human gold, and neither may replace the blank full-cohort or full-impact real-person packets.

The reviewer-C tiebreak reduces missing candidate coverage but adds post-selection dependence. Its worklist is blind to A/B content and its sessions are disjoint, yet the 103 rows are selected only after A/B unsealing and all three roles share one model family, prompt, and snapshot. The 66 resolved rows are therefore an opportunistic same-family majority candidate, not a preregistered independent third expert. The fixed gate correctly remains no-go at 0.6408 selected-row resolution and 0.9704 combined coverage. Reporting baseline agreement on the 1,213 selected candidates does not estimate accuracy, and the unresolved 37 rows may be systematically harder than the retained rows.

The D/E evidence-secondary audit is selected after both A/B and C are unsealed and therefore cannot repair that selection dependence. Its URL selector is frozen before new retrieval and both reviewer inputs are blind to prior labels, but all five roles still share one model family and the 37-row cohort is explicitly the hardest remainder. Nine non-human rows receive no successful non-empty evidence record, while successful page retrieval does not imply complete field evidence. Exact non-human D/E agreement on 25/37 is dominated by 19 shared uncertain decisions; only 6 non-human rows are evidence-qualified strict, and non-human affected\_versions resolution remains 0/28. These values characterize abstention and evidence limitations in a selected cohort, not independent-expert agreement, prevalence, or human-gold accuracy.

The residual non-affected audit is even more targeted: it selects the final two CWE rows and one references row after D/E outcomes are known, and its source, taxonomy, and URL profile shapes are inspected before v1. The 1/3 mechanism-supported result therefore cannot estimate residual resolution probability or validate the \texttt{factual\_conflict} candidate. Absence of a local \texttt{try} supports an uncaught-exception path but does not prove that resource consumption is impossible elsewhere. Likewise, a prohibited CWE category does not prove the alternative source assignment, and a malformed-URL repair is a construct choice rather than an observed redirect. All three promoted fields remain null and combined coverage is unchanged.

The staged frontier is also post-hoc and conditions each marginal ratio on the preceding stage's failures. Its 66/103 and 6/37 yields describe different, progressively harder selected sets and must not be read as a controlled model comparison or extrapolated to real reviewers. Token totals cover successful responses only because the three missing B attempts have no usage record. The stop rule is valid only for further same-model escalation on this revealed cohort; it does not establish that future models, real reviewers, or an independent snapshot would have the same yield.

The non-human diagnostic holdout also reveals an evaluator-contract threat. The current severity baseline uses canonical labels, whereas the reviewer protocol defines the field over labels, scores, vectors, and CVSS versions; 165 strict disagreements therefore combine comparator error with a deliberate difference in construct. The affected\_versions projection separately removes package-specific unbounded records whose only range marker is \texttt{introduced=0}, producing 25 strict one-sided-presence disagreements. A post-hoc profile that encodes these observed mechanisms reaches 1,121/1,147 on the same non-human strict rows, but this is in-sample diagnostic fit and cannot estimate generalization. The blank full-cohort human packet reduces selection bias in the planned audit, yet it has 0/1,250 signed rows and its validator cannot prove that supplied identities correspond to independent real people; identity and review independence require external verification.

This non-human diagnostic shows that protocol-generation drift prevents retrospective pooling. Replaying the fresh severity contract on the older AI-gold rows reverses direction, but both generations still use related Codex workflows, so the reversal does not identify which contract is correct. For affected\_versions, the old seed stores empty projected range values for ten one-sided-package rows; current raw aligned data cannot be substituted as the label-time input. Consequently, neither a pooled accuracy estimate nor current-snapshot reconstruction is valid evidence for the mechanism. A shared real-human calibration under one explicit contract must precede any new candidate seal.

The two AI contract calibrations reduce input and ordering variation but do not remove model-family or selection dependence. Both draw from the already reviewed fresh cohort and require prior non-human strict consensus for selection. V2 is row-disjoint from v1, but its refined clauses were derived after v1 unsealing, so it validates reviewer stability under the revised wording rather than external correctness. The conditional XWiki evidence secondary and both deterministic projection passes are selected after v2 unsealing. Projection v2 closes the earlier coordinate gap with official product POMs and source continuity, but the chosen class set, edge schema, and release parser were developed on this known failure. Its \texttt{incomplete} candidate is therefore a traceable case explanation, not an independent estimate, human agreement, or gold accuracy. The subsequent generic graph contract is fixed before cross-case evidence retrieval, but both cross-case sources had already been reviewed and derive from a non-human-label-conditioned calibration. Equality selection makes the first 8/8 result an expected construct-consistency diagnostic. The non-equal selector avoids reviewer files and excludes XWiki, but its 2/5 agreement diagnostic still compares against same-model-family labels rather than human truth. Registry catalogs additionally implement extensional snapshot semantics and can omit prereleases, as Graylog demonstrates. The later Codex contract candidate explicitly chooses that extensional interpretation after the no-go is known and tests only one InLong row whose two component ranges are identical. The three-row unseen-ecosystem cohort improves selection independence because it comes from the full aligned input rather than consensus-conditioned calibration data, but its ecosystems and structural filters were chosen after the earlier no-go, it contains one row per ecosystem, and all initial product projections abstain. The Deno lockfile recovery contract is fixed before its new evidence fetch, yet it is explicitly designed after inspecting that row's dependency-constraint failure and reuses the same revealed CVE. Its 71 complete mappings support a one-case mechanism recovery, not prevalence, accuracy, or an unbiased success rate. Generalization still requires a human-approved range contract, edge-class mapping rules tested beyond one project, a later untouched cohort, and real-person validation.

The 28-row edge audit is also post-unsealing. Its family rules and structural score allocate development effort; they do not estimate a family success probability. Mattermost then undergoes two explicit protocol revisions before v3 set analysis: v1 exposes an API pagination assumption, and v2 shows that supplied version boundaries need not be GitHub Release objects. V3 narrows the construct to 19 exact Git tags after those failures are known. The cache-only verifier establishes the reported manifest and ancestry facts, but the 0/2 result remains a selected two-row mechanism diagnostic. It cannot be interpreted as Mattermost ground truth, a Go-ecosystem failure rate, or an unbiased evaluation of graph projection.

LF Edge EVE is likewise selected after the 28-row audit and undergoes disclosed protocol discovery before v1. The 207-tag grammar, absent root module, nested component manifests, and pseudo-commit changed paths are known when the contract is frozen, so v1 explicitly disables candidate promotion. The Git pack and independent verifier establish the reported ref, object, timestamp, path, and ancestry facts, but the component and patched-anchor gates compare a public advisory API snapshot with detail observed in a different repository-advisory interface. Interface drift or representation differences may explain that mismatch. The 0/2 result is therefore a selected mechanism and evidence-interface diagnostic, not EVE ground truth, a historical-advisory correction, or a Go-repository prevalence estimate.

Hutool is also selected from the revealed 28-row audit, and its equal 214-token catalogs, 209-token stable grammar, and three aggregate anchors are inspected before v1 is frozen. The independent verifier establishes the cached catalog, POM, JAR, and set facts, but source/JAR containment is inspected at three critical anchors rather than all 209 releases. More importantly, the projection explicitly chooses current-snapshot extensional semantics for two GHSA ranges with no upper bound. The 2/2 mechanism pass and two development \texttt{incomplete} candidates therefore cannot be interpreted as corrected labels, candidate coverage, historical advisory intent, Maven-family prevalence, or human agreement. They only justify evaluating the pre-frozen mechanism on a later blind cohort.

The six-row Hutool external application is disjoint by CVE from the 1,967-record prior-exposure union, and neither baseline nor reviewer labels enter selection. Nevertheless, it is not an untouched temporal holdout: the records come from the same aligned snapshot and their availability and structural routes were inspected before sealing. Its 6/6 projection result therefore supports only retrospective mechanism reuse under the declared snapshot-extensional contract. It cannot estimate future-snapshot coverage, Maven-family prevalence, taxonomic accuracy, or human agreement, and none of its six candidates is promoted.

The leave-one-cohort-out analysis also remains post-hoc: its authority classes, structured features, models, and thresholds were selected after v2 unsealing, and all three cohorts use non-human labels from related workflows. CVE-disjoint splits reduce row leakage but do not remove shared candidate-miner, evidence-source, or model-family dependence. The pooled type gain is not stable on held-out Phase D, and the source models do not exceed the strongest named baseline. We use this analysis only as a no-go gate, not as confirmatory model selection.

The v2 reviewer files also required disclosed mechanical repairs by their original reviewers: one JSON null was converted to an empty string, and short source rationales were expanded without changing labels, sources, or evidence. The frozen merge contract's character limits were not stated explicitly in the reviewer prompt. The repairs do not alter the substantive decisions, but the prompt/validator mismatch is a protocol defect and weakens any claim of a fully hands-off annotation run.

Reproducibility depends on preserving generated artifacts and manifests. The paper points to generated JSON, Markdown, CSV, and SVG outputs so counts can be rechecked, but final publication still needs a stable dependency list, final bibliography format, and a clear statement of which data snapshots were used.

\section{Related Work}

This paper sits between four bodies of prior work.

First are studies and source documents that document inconsistency in vulnerability data sources. VIEM \cite{VIEM2019} shows that extracting affected software and versions from unstructured reports is hard and that agreement with NVD is far from complete. Croft et al. \cite{Croft2022} show that severity differs across sources and that those differences can affect downstream prediction. Croft et al. are the closest structured severity study, but their scope is one project/source family and they do not type non-conflict differences or adjudicate source support. NVD-focused audits such as "Cleaning the NVD" \cite{Anwar2022} and Zhang et al.'s \emph{The Flaw Within} \cite{Zhang2023} show that even a single database is not immune to quality problems, while user studies of NVD usage and CVSS scoring inconsistency show that practitioners encounter missing or incorrect entries and that severity scoring can be subjective even among evaluators \cite{NVDUsers2024,Wunder2024}. NVD process documents and recent operations updates further clarify that enrichment is an active analysis process, not a mechanical copy of CVE records, and that enrichment priorities can change under record-volume pressure \cite{NVDProcess2026,NVDStatus2026,NVDTransition2024,NVDOperations2026}. Recent affected-version benchmarking further shows that identifying affected versions is difficult even when tools use commits and advisories \cite{Chen2025AffectedVersions}. These works motivate the problem, but they do not directly solve post-alignment discrepancy typing for structured NVD-GHSA field pairs or evidence-constrained adjudication of residual factual conflicts. Unlike schema-cleaning or single-database repair work, this paper does not overwrite records or produce a canonical cleaned vulnerability database; it preserves both source values and emits a discrepancy type plus routing action.

Second are papers that detect inconsistency in unstructured vulnerability text. Sun et al.'s TOSEM work on aspect-level discrepancies \cite{Sun2023}, Wang et al.'s aspect extraction work on heterogeneous threat intelligence \cite{Wang2025Aspects}, and Li et al.'s 2025 \emph{VuldiffFinder} \cite{Li2025VuldiffFinder} are the closest neighbors in spirit, but they operate on report text and aspect extraction, not on already aligned structured NVD-GHSA fields. GapFinder \cite{GapFinder2021} is similar in method family but is aimed at broader security-intelligence text rather than aligned vulnerability records. VuldiffFinder is especially important because it is a recent Computers \& Security paper on vulnerability inconsistency; the boundary is that it discovers inconsistencies in unstructured vulnerability information, while this paper starts after database normalization/CVE alignment and studies type-first structured-field discrepancy handling. In this setting, the method output is not just an inconsistency flag: structured fields introduce schema-specific non-conflict mismatch modes such as subset reference lists, same-day timestamp representations, and range-bound shape differences, so the workflow emits an operational discrepancy type and routing action before any source-support question is asked.

Third are generic truth-discovery and multi-source fusion frameworks. CRH \cite{CRH2016} and truth-discovery surveys \cite{Wang2024TruthSurvey} provide methodological background for resolving conflicts across heterogeneous sources. Standard truth-discovery frames conflict resolution as true-value/source-reliability estimation; this setting instead uses explicit advisory evidence, field semantics, and an abstain/manual-review outcome. A CRH-style global reliability model is not the right first operation here because the main comparison has two source families and many post-alignment mismatches are representation, incompleteness, or timing differences rather than competing truth claims. Selective classification work \cite{Goren2024HSC} is relevant because it formalizes abstention under uncertainty, although this paper uses evidence and vulnerability-field semantics rather than learned classifier confidence.

Finally, vulnerability-intelligence ecosystem studies explain why NVD-GHSA discrepancies matter operationally. GHSA documentation and repository metadata show that the advisory database combines GitHub-reviewed advisories, unreviewed submissions, NVD imports, OSV-format records, and GitHub-specific affected-version extensions \cite{GHSAAbout2026,GHSARepo2026}. Work on the GHSA review pipeline highlights timing differences in advisory processing \cite{Segal2026GHSA}, VulZoo illustrates the scale of multi-source vulnerability intelligence integration \cite{Ruan2024VulZoo}, and scanner/evaluation studies show that heterogeneous sources, identifier schemes, and ambiguous version ranges can affect downstream vulnerability detection results \cite{Churakova2026VEX,Mandl2026VulnDetection}. These works do not solve field-level adjudication, but they motivate temporal discrepancies, source coverage gaps, affected-version ambiguity, and downstream scanner divergence as practical concerns.

The table below summarizes the boundary in reviewer-facing terms. The point is not that structured fields are easier than text; rather, structured fields fail in different ways after alignment, so a method needs field semantics, discrepancy types, and an abstention-capable evidence handoff.

\begin{table}[!htbp]
\centering
\setlength{\tabcolsep}{3pt}
\scriptsize
\caption{Reviewer-facing boundary between this paper and neighboring lines of work. The current contribution is a post-alignment structured-field typing workflow with an evidence-constrained adjudication scaffold.}
\label{tab:08-related-table-1}
\begin{tabularx}{\textwidth}{>{\raggedright\arraybackslash}p{0.18\textwidth}>{\raggedright\arraybackslash}X>{\raggedright\arraybackslash}X>{\raggedright\arraybackslash}X>{\raggedright\arraybackslash}X}
\toprule
Line of work & Input and stage & Output usually studied & Evidence and abstention & Boundary relative to this paper \\
\midrule
Vulnerability database quality, schema cleaning, and inconsistency studies & Single-source audits, schema repair, or cross-source observations & Missingness, quality defects, canonicalized records, severity differences, affected-version difficulty & Usually descriptive or repair-oriented; source support is not the main output & Motivates the problem, but this paper keeps both source values and types post-alignment field discrepancies instead of producing a canonical clean database \\
Unstructured advisory or threat-intelligence discrepancy mining & Report text, advisory prose, extracted aspects & Textual or aspect-level inconsistency labels & Evidence is the input text; abstention is usually not a first-class outcome & Closest in spirit, but this paper starts after CVE alignment and outputs a structured-field type plus routing action before source-support adjudication \\
Generic truth discovery and data fusion & Multi-source claims before domain-specific field routing & True-value and source-reliability estimates & Evidence is often implicit in source agreement; abstention is not central & This paper does not learn a global truth model; it first removes non-conflict field mismatches, then uses fetched advisory evidence and abstention for residual conflicts \\
Affected-version and scanner-ecosystem work & Package identities, range events, commits, scanner outputs & Applicability errors, scanner divergence, version-range difficulty & Evidence may include commits and advisories, but not this field-typing workflow & Explains the hardest field; the current method remains diagnostic until package-aware and range-aware adjudication is validated \\
\bottomrule
\end{tabularx}
\end{table}

Our paper is narrower and more concrete: after CVE-ID alignment, it types field-level discrepancies on structured NVD-GHSA pairs, then prototypes evidence-constrained adjudication with abstention on sampled factual conflicts. The current results are not a replacement for full truth discovery and are not a claim that silver labels are gold. The defensible contribution boundary is this type-first, post-alignment workflow: structured-field discrepancy typing first, and evidence-constrained source-support analysis only for residual factual conflicts.

\section{Conclusion}

This paper studies what remains after CVE-ID alignment: structured fields can still disagree in ways that require different downstream actions. The aligned NVD-GHSA corpus shows substantial field-level discrepancy after matching, and the RQ2 diagnostics show that most binary differences are not routed to deterministic factual-conflict candidates. A 1,250-CVE exposure-disjoint typing-stability cohort reaches 0.9176 strict same-model consensus coverage. A post-unseal third blind pass raises non-human candidate coverage to 0.9704 but fails its gate. A subsequent frozen-URL D/E audit resolves only 6 of the remaining 37 non-human rows and raises coverage to 0.9752, again below fixed thresholds; affected\_versions remains 0/28. A targeted audit of the remaining two CWE rows and one references row supports one additional development mechanism candidate but promotes none because the cohort and evidence routes are already revealed. An immutable-log frontier reconstructs 2,767 request row-attempts, including 90 retry row-items, and shows marginal selected-row yield falling from 66/103 at C to 6/37 at D/E; the fixed stop rule rejects further same-model escalation on this cohort. A later official snapshot contains no strict event-time CVE but supports a 250-row snapshot-external development cohort with 0.9240 strict dual-Codex coverage. Only three rows distinguish sealed profiles, and only one has strict consensus favoring the CWE candidate in the main review. A post-selection three-case frozen-evidence diagnostic makes two fresh same-family passes agree with the candidate direction on all three cases. Expanding the same diagnostic to all 50 CWE rows retains 49 strict decisions and yields 45/49 current versus 48/49 candidate agreement, but field selection remains post-hoc. A final 16-row non-CWE evidence stage resolves only four rows despite complete fetch availability; its staged 238/250 candidate passes coverage but fails the fixed 0.40 resolution gate and leaves the sealed evaluation unchanged. A label-free enumeration proves that the three profile-difference rows bound any full-cohort accuracy difference to +/-1.2 percentage points under every possible five-label assignment. All six profiles collapse to two vectors on that sample, and no possible label assignment can make their exact two-sided paired comparison significant at alpha 0.05; the minimum attainable p-value is 0.25. A prediction-only census over all 5,948 eligible CVEs instead separates all six vectors but finds only 34 difference CVEs across 29,740 field instances: 29 CWE and at most five references differences. This revealed-universe result corrects the sampled-rate assumption and informs disagreement-enriched future design; without labels it does not show correctness or gain. None of the evidence stages permits promotion. Complete-impact references and CWE audits expose additional bounded candidate directions and construct disagreements. The affected\_versions holdouts separate discrepancy typing from source adjudication and find no stable method improvement. All current annotation-backed comparisons remain non-human rather than human-gold performance results.

The main methodological lesson is type-first: post-alignment fusion needs a discrepancy type before it needs a winner. Many field differences are better handled as normalization, incompleteness, or temporal interpretation, and only residual factual conflicts should be passed to evidence-constrained source-support analysis.

The present contribution is therefore a scoped task definition, taxonomy, auditable baseline, evidence-handoff scaffold, and leakage-controlled, task-separated evaluation protocol with explicit failure, protocol-drift, evidence-dependence, event-time eligibility, and prediction-identifiability analysis. It is not validated as an adjudication method: the v2 type candidate is too selective, the source method does not beat a fixed-source comparator, and the RQ2 post-hoc contract profile reverses direction across protocol generations. Two disjoint AI calibrations refine the candidate contract and stabilize four clauses, but the full gate still rejects it. A targeted residual audit shows that frozen source code can support one local CWE mechanism while official mapping validity and URL-resource identity leave two cases unresolved; its one development \texttt{factual\_conflict} is not promoted. For the remaining XWiki row, a lineage-aware graph binds product dependencies and artifact migration and derives an \texttt{incomplete} development candidate with one missing prerelease. A subsequent generic contract binds eight equal-boundary artifact-mismatch cases, but a fixed non-equal test reaches only 4/5 projection coverage and 2/5 consistency with both sealed AI reviewers. A frozen snapshot-extensional Codex candidate projects two homogeneous InLong components but disagrees with both reviewers. A label-independent full-aligned-data cohort across NuGet, PyPI, and crates.io initially reaches 0/3 product projections because package catalogs and dependency constraints lack total product-release mappings. A fixed post-no-go Deno diagnostic then recovers one such edge from 71 official tag lockfiles and derives an NVD-strict-subset, non-human \texttt{incomplete} candidate with five runtime-only releases. By contrast, Mattermost and EVE each yield 0/2 project-family projections, while the Hutool mechanism projects two development rows and six same-snapshot exposure-disjoint rows without promotion. Together these cases show why edge-specific build evidence can recover one sequence while repository identity, valid pseudo syntax, commit timestamps, or current catalog equality alone cannot establish an approved historical range construct. Completing validation still requires a human-approved range contract and edge-class-specific mapping evidence, followed by two real reviewers and author signoff over both prepared 1,250-row and post-profile 250-row blind RQ2 packets and the RQ3 human audit. A future confirmation cohort must require both NVD and reviewed-GHSA event timestamps after the frozen profile, census frozen predictions before label collection, and separate disagreement-focused paired testing from probability-sampled absolute evaluation; a later download alone is insufficient.


\input{sections/declarations.tex}

\bibliographystyle{elsarticle-num}
\bibliography{references}

\end{document}
